\documentclass[11pt]{article}
\usepackage[margin=1in]{geometry}
\usepackage[T1]{fontenc}
\usepackage[utf8]{inputenc}
\usepackage{lmodern,amsmath,amssymb,amsthm,mathtools,booktabs,microtype}
\usepackage[colorlinks=true,linkcolor=blue,citecolor=blue,urlcolor=blue]{hyperref}
\hypersetup{pdftitle={Finite Gaussian Bayesian networks can compete with small approximating networks},pdfauthor={ChatGPT-6 Astra; Claude Opus 5.5; Kaarel H\"anni; Dmitry Vaintrob}}
\newtheorem{theorem}{Theorem}[section]
\newtheorem{lemma}[theorem]{Lemma}
\newtheorem{proposition}[theorem]{Proposition}
\newtheorem{corollary}[theorem]{Corollary}
\theoremstyle{remark}\newtheorem{remark}[theorem]{Remark}
\newcommand{\E}{\mathbb E}\newcommand{\R}{\mathbb R}\newcommand{\N}{\mathbb N}
\newcommand{\Prob}{\mathbb P}\newcommand{\clip}{\operatorname{clip}}
\newcommand{\KL}{\operatorname{KL}}\newcommand{\Var}{\operatorname{Var}}
\newcommand{\Rad}{\operatorname{Rad}}
\title{Finite Gaussian Bayesian networks can compete with small approximating networks}
\author{ChatGPT-6 Astra \and Claude Opus 5.5 \and Kaarel H\"anni \and
Dmitry Vaintrob\thanks{Authors are listed in alphabetical order.}}
\date{17 November 2026}
\begin{document}
\maketitle
\begin{abstract}
We prove squared-loss comparisons with the best small approximating network for finite Bayesian posteriors of fully connected tanh networks. Inputs and labels may have any bounded IID law; no exactly realizable teacher or noise model is assumed. The predictor is the ordinary posterior mean of the raw network output. At every learner width above a competitor's width, its risk is at most the competitor's risk plus a term that vanishes polynomially with sample count. The learner has independent Gaussian weights with the fixed width scaling $\sigma^2=N^{-\gamma}$, for every fixed $\gamma\ge1$, and a calibrated Gaussian likelihood. For $H\ge2$ hidden layers, the remainder at fixed competitor resources is $\widetilde O(n^{-1/H})$ when $\gamma\ge2$, and $\widetilde O(n^{-1/(2^{H+1}+3)})$ at mean-field scaling $\gamma=1$. A common polynomial calibration covers all exponents and widths on one sample event. One hidden layer gives a square-root rate for all $\gamma\ge1$. A sample-dependent overall prior prefactor recovers the sharper deep-network rate throughout this range. The proof combines norm-based capacity, Gaussian spectral truncation, and finite covers built from bounded-relative-entropy covariance sets. No prior over network sizes is introduced.
\end{abstract}

\section{Competing with an approximating network}
Many prediction problems have no small network that describes the data exactly. A more useful comparison asks whether a Bayesian neural network, after $n$ observations, can approach the best population loss attainable by a much smaller network. Write
\[
 R(g)=\mathbb E(g(X)-Y)^2.
\]
We prove bounds of the form
\[
 R(\widehat f_S)\leq\inf_{g\in\mathcal G_n}R(g)+o(1),
\]
where $\mathcal G_n$ contains networks of polynomially growing width and bounded weights, and the learner may be arbitrarily wider. The leading coefficient on the comparison loss is one. Both model misspecification and bounded label noise are allowed.

An additive remainder is necessary even for estimating an unknown constant mean. A weight bound is also material: a narrow network with large weights can vary on arbitrarily small input scales, while a Gaussian prior charges those weights. Our comparison classes therefore specify width and a weight cap. Appendix~\ref{width-only:sec:width-only-obstruction} makes this obstruction precise even in one dimension, with noiseless labels and input variance bounded below. These comparison classes are used only to state the guarantee; the learner is one fully connected network with a Gaussian prior, not a mixture over comparator architectures.

\paragraph{The learner.}
There are $H$ hidden tanh layers, a scalar linear readout, a constant first-layer input feature, and no internal biases. The total depth is $L=H+1$. At learner width $N$, the Gaussian weight variances are
\[
 \frac{N^{1-\gamma}}d\quad\hbox{in the first layer},\qquad
 N^{-\gamma}\quad\hbox{in internal layers},\qquad
 N^{-\gamma-1}\quad\hbox{in the readout}.
\]
For a deterministic calibration $\lambda_n>0$, the observation variance is $n\lambda_nN^{-\gamma}$. This is a tempered likelihood: it balances fit against the prior rather than requiring interpolation of noisy labels. Our results concern exact posterior inference; they do not provide an efficient inference algorithm.

\paragraph{Finite priors at every exponent $\gamma\ge1$.}
For any fixed input dimension, hidden depth, teacher resources, and confidence level, the following polynomial remainder bounds hold. All logarithmic factors are suppressed in this introductory table.
\begin{center}
\begin{tabular}{lll}
\toprule
Hidden layers & Width exponent & Remainder rate \\
\midrule
$H=1$ & every $\gamma\ge1$ & $n^{-1/2}$ \\
$H\ge2$ & every $\gamma\ge2$ & $n^{-1/H}$ \\
$H\ge2$ & $\gamma=1$ & $n^{-1/(2^{H+1}+3)}$ \\
$H\ge2$ & each $1<\gamma<2$ & $n^{-\kappa_{H,\gamma}}$ \\
\bottomrule
\end{tabular}
\end{center}
Here
\[
 b_H=\max\{1,(H-1)/2\},\qquad
 \kappa_{H,\gamma}=
 \frac{\gamma-1}{3(\gamma-1)+2b_H\gamma}>0.
\]
There is also one calibration, independent of $\gamma$, with remainder
\[
 \widetilde O\!\left(n^{-1/(2^{H+1}+2b_H+3)}\right)
\]
on a single sample event for every $\gamma\ge1$ and every integer learner width above the competitor's width (Theorem~\ref{thm:fixed-all-gamma-raw}). This common rate is $n^{-1/13}$, $n^{-1/21}$, and $n^{-1/38}$ for two, three, and four hidden layers. It improves the intermediate-scale guarantee near $\gamma=1$; the sharper separate endpoint bounds remain useful. No width limit or prior-strength limit is taken.

\paragraph{How large may the competitor be?}
Let $m$ be its width and $B\ge1$ its weight cap. A useful implementation-energy bound is
\[
 E=d(d+1)B^2+2Hm^2B^2.
\]
For $H\ge2$ and $\gamma\ge2$, width and weight cap may both be $n^\alpha$ for any $\alpha<1/(4H)$. At a fixed weight cap, width may be $n^\alpha$ for any $\alpha<1/(2H)$. At mean-field scaling the corresponding ranges are $\alpha<2^{-(H+3)}$ and $\alpha<2^{-(H+2)}$. A single all-exponent calibration permits both resources to grow as $n^\alpha$ for any $\alpha<1/[4(2^{H+1}+2b_H)]$, or width alone at twice this exponent when the weight cap is fixed. Every fixed intermediate exponent has the additional, sometimes sharper, range in Section~\ref{sec:fixed-above-one-raw}. For two hidden layers the common all-exponent ranges improve to $\alpha<1/36$ and $\alpha<1/18$, respectively. The calibrations depend on the specified class caps, not on the unknown data law or its best approximant.

With one hidden layer, the ranges improve to $\alpha<1/4$ when width and weight cap grow together and $\alpha<1/2$ when the weight cap is fixed. These are comparisons with the best member of a compact class: a minimizing teacher exists and can be chosen before the sample is drawn.

\paragraph{An additional scale parameter.}
Replacing $N^{-\gamma}$ by $t_n^2N^{-\gamma}$, with an explicit positive $t_n$ decreasing with sample count, gives the faster $\widetilde O(n^{-1/H})$ deep-network rate for every $\gamma\ge1$. This remains a finite prior and one fixed architecture at each width. It is, however, a different sequence of priors indexed by sample count; the theorem is stated separately from the fixed-prefactor results.

\paragraph{Why learner parameter counting does not give these bounds.}
The learner's parameter count may greatly exceed $n$. The proof instead charges a functional complexity to the prior. For stronger scales, weighted norm and smooth-loss bounds suffice. In the intermediate regime, a typical Gaussian internal matrix has large Frobenius norm but much smaller operator norm; separating these contributions gives a more useful effective energy. At mean-field scaling, we cover the outputs associated with bounded relative entropy at each covariance layer. A Gaussian auxiliary measure on one finite architecture counts these approximants; it is solely a proof device. Posterior tail estimates then extend the bounds to raw, potentially unbounded predictions.

\section{Model, risk, and prediction convention}
\label{sec:gaussian-oracle}
Let $d\ge1$, $H\ge1$, and $N\ge1$. Inputs lie in $[-1,1]^d$ and labels in $[-1,1]$. Their joint Borel law $D$ is otherwise arbitrary. A sample $S=((X_i,Y_i))_{i=1}^n$ is drawn independently from $D$. Set $\bar x=(1,x)$ and
\[
 h^1(x)=\tanh(W^1\bar x),\qquad
 h^\ell(x)=\tanh(W^\ell h^{\ell-1}(x))\quad(2\le\ell\le H),
 \qquad z_w(x)=v^\top h^H(x).
\]
There are no internal biases in the learner. This does not restrict the comparison to teachers without biases: a width-$m$ tanh teacher with all weights and hidden/readout biases bounded in absolute value by $B\ge1$ compiles exactly into this architecture at width $m+1$ and weight cap $2B$. Carry one constant neuron of value $1/2$ through each layer. Its first preactivation is $\operatorname{arctanh}(1/2)$ and each later self-edge is $2\operatorname{arctanh}(1/2)=\log3$; a bias $b$ becomes an incoming edge of weight $2b$ from this neuron. The same construction handles a readout bias. This comparison construction does not change the learner's prior. The total depth, including the linear readout, is $L=H+1$. The weighted energy is
\begin{equation}
 Q_N(w)=\frac dN\|W^1\|_F^2+\sum_{\ell=2}^H\|W^\ell\|_F^2+N\|v\|_2^2.
 \label{eq:gaussian-energy}
\end{equation}
For statistical bounds we also use the energy after the first layer,
\[ Q_N^+(w)=\sum_{\ell=2}^H\|W^\ell\|_F^2+N\|v\|_2^2. \]
The prior $\mu_{N,\sigma}$ gives all weight coordinates independent centered Gaussian laws with variances $N\sigma^2/d$, $\sigma^2$, and $\sigma^2/N$ in the three layer types. Its density is proportional to $\exp[-Q_N/(2\sigma^2)]$. Given $\lambda>0$, observations have Gaussian likelihood with mean $z_w(X_i)$ and variance $n\lambda\sigma^2$. Thus
\[
 \Pi_{N,\sigma,S}(dw)=Z^{-1}
       \exp\!\left[-\frac{\widehat R_S(z_w)}{2\lambda\sigma^2}\right]
       \mu_{N,\sigma}(dw),
 \qquad \widehat R_S(f)=\frac1n\sum_i(f(X_i)-Y_i)^2.
\]
Our main predictor is $\widehat f_S(x)=\int z_w(x)\,d\Pi_{N,\sigma,S}(w)$.
The stronger bounds proved below control $\int R(z_w)\,d\Pi$, so Jensen's inequality gives the same bound for $R(\widehat f_S)$. The readout and bounded hidden activations imply $|z_w(x)|\le\sqrt{Q_N(w)}$; all these integrals are finite. Clipping to $[-1,1]$ can only decrease squared loss against these labels, but is not part of the main prediction convention.

\section{Teacher implementations and a width-independent capacity bound}

\begin{lemma}[Cloning a teacher]
\label{lem:gaussian-cloning}
Suppose that $g$ is a network of the preceding architecture, hidden
width $m$, and depth $H+1$, with every weight of absolute value at most
$B$.  At every width $N\geq m$ it has an exact implementation with
\begin{equation}
 Q_N\leq E:=d(d+1)B^2+2H m^2B^2.
 \label{eq:gaussian-teacher-energy}
\end{equation}
Moreover, $|g(x)|\leq U:=mB$ for every $x$.
\end{lemma}
\begin{proof}
Partition the $N$ neurons in each hidden layer into $m$ groups, of sizes
$k_j\in\{\lfloor N/m\rfloor,\lceil N/m\rceil\}$.  Copy a teacher
first-layer row onto every neuron in its group.  Replace a teacher
internal weight $A^\ell_{ji}$ by $A^\ell_{ji}/k_i$ on every corresponding
edge, and its readout weight $u_j$ by $u_j/k_j$ on each copy.  Every copy
has exactly the teacher activation, so the function is unchanged.
The first-layer energy is at most $d(d+1)B^2$.  Since $k_j/k_i\leq2$,
each internal layer costs at most $2m^2B^2$.  Since
$k_j\geq N/(2m)$, the readout costs at most $2m^2B^2$.  There are
$H-1$ internal layers and one readout.  Finally, $|\tanh|\leq1$ gives
$|g|\leq\sum_j|u_j|\leq mB$.
\end{proof}

Put
\begin{equation}
 D_n=\sqrt{2(d+1)\log(n+1)}+\frac{H-1}{2},\qquad
 \mathcal Z_K=\bigcup_{N\geq1}\{z_w:Q_N^+(w)\leq K\}.
 \label{eq:gaussian-Dn}
\end{equation}
The union over all learner widths is essential: the following event
controls every width at once.

\begin{lemma}[Capacity of an energy ball]
\label{lem:gaussian-capacity}
For every set of $n$ inputs and every $K>0$, the empirical Rademacher
complexity satisfies
\begin{equation}
 \widehat{\mathfrak R}_S(\mathcal Z_K)
 \leq \left(\frac KH\right)^{H/2}\frac{D_n}{\sqrt n}.
 \label{eq:gaussian-rademacher}
\end{equation}
Consequently, with probability at least $1-\eta$, simultaneously for all
widths and all $w$ with $Q_N^+(w)\leq K$,
\begin{equation}
 R(\operatorname{clip}z_w)
 \leq\widehat R_S(\operatorname{clip}z_w)
 +8\left(\frac KH\right)^{H/2}\frac{D_n}{\sqrt n}
 +4\sqrt{\frac{\log(1/\eta)}{2n}}.
 \label{eq:gaussian-fixed-energy-deviation}
\end{equation}
\end{lemma}
\begin{proof}
For a bounded set $A\subset\mathbb R^n$, write
$\operatorname{Rad}(A)=n^{-1}\mathbb E_\epsilon
\sup_{a\in A}\sum_i\epsilon_i a_i$.
Its value is unchanged by taking a closed convex hull; it scales
linearly under positive scalar multiplication.  We use the standard
Rademacher contraction inequality, Massart's finite-class bound, and
the Sauer--Shelah bound for halfspaces.

For $t\in[0,1]$, let $\operatorname{clip}_t(u)=\max\{-t,\min\{t,u\}\}$,
and let $\operatorname{Clip}(A)$ contain the vectors
$(\operatorname{clip}_t(a_i))_{i=1}^n$ for $a\in A$ and a single common
$t\in[0,1]$.  The identity
$\operatorname{clip}_t(u)=(|u+t|-|u-t|)/2$, contraction, and symmetry of
the Rademacher signs imply
\begin{equation}
 \operatorname{Rad}(\operatorname{Clip}(A))
 \leq\operatorname{Rad}(A)+\frac1{2\sqrt n}.
 \label{eq:gaussian-variable-clip}
\end{equation}
Indeed, split the supremum between the two absolute-value expressions
and contract each.  The two additional terms sum to
$\mathbb E|\sum_i\epsilon_i|/(2n)\leq1/(2\sqrt n)$.

We also need a positive mixture identity.  For $s>0$ and $|u|\leq1$,
\begin{equation}
 \frac{\tanh(su)}s=
 \int_0^1\operatorname{clip}_t(u)\,
       2s\tanh(st)\operatorname{sech}^2(st)\,dt
 +\operatorname{sech}^2(s)\operatorname{clip}_1(u).
 \label{eq:gaussian-tanh-mixture}
\end{equation}
All weights are nonnegative and have total mass one.  To verify the
identity for $u\geq0$, integrate by parts with
$\psi(t)=\tanh^2(st)$: the right side is
$u-\int_0^u\tanh^2(st)dt=\tanh(su)/s$.  Oddness gives the other half.

On the fixed sample inputs, let $\mathcal F_1$ be the closed convex hull
of the vectors $(\tanh(a\cdot\bar X_i))_{i=1}^n$, with
$a\in\mathbb R^{d+1}$.  Recursively set
$\mathcal F_\ell=\overline{\operatorname{conv}}
\operatorname{Clip}(\mathcal F_{\ell-1})$.  These sets are symmetric,
convex, contain zero, and lie in $[-1,1]^n$.
Define $\rho^1=\boldsymbol1_N$ and
$\rho^\ell=|W^\ell|\rho^{\ell-1}$, with entrywise absolute values.
Induction using \eqref{eq:gaussian-tanh-mixture} shows that the sample
vector of hidden neuron $j$ at layer $\ell$ belongs to
$\rho^\ell_j\mathcal F_\ell$.  A zero value of $\rho^\ell_j$ forces the
neuron to vanish.  For a positive value, divide the preactivation by
$\rho^\ell_j$; it is a convex signed combination of the previous
normalized neurons, hence belongs to $\mathcal F_{\ell-1}$.  The mixture
identity completes the induction.

It follows that the output sample vector belongs to
$P(w)\mathcal F_H$, where
$P(w)=|v|^\top|W^H|\cdots|W^2|\boldsymbol1_N$ (with no matrix factors
when $H=1$).  Cauchy--Schwarz and the Frobenius bound on operator norm give
\[
 P(w)\leq\sqrt N\|v\|_2\prod_{\ell=2}^H\|W^\ell\|_F
 \leq\left(\frac{Q_N^+(w)}H\right)^{H/2}.
\]
The last inequality is AM--GM over the $H$ energy terms after the first
layer.  There is no restriction on first-layer weights in this capacity estimate. The positive quadratic form on the internal and readout coordinates suffices.
By \eqref{eq:gaussian-variable-clip},
$\operatorname{Rad}(\mathcal F_H)\leq
\operatorname{Rad}(\mathcal F_1)+(H-1)/(2\sqrt n)$.

Finally, if a scalar random variable $T$ has distribution function
$(1+\tanh t)/2$, then $\tanh u=\mathbb E\operatorname{sign}(u-T)$.
Thus $\mathcal F_1$ lies in the convex hull of the affine-halfspace
dichotomies of the sample inputs.  There are at most $(n+1)^{d+1}$ such
dichotomies by Sauer--Shelah.  Massart's bound gives
$\operatorname{Rad}(\mathcal F_1)
\leq\sqrt{2(d+1)\log(n+1)/n}$, proving
\eqref{eq:gaussian-rademacher}.

The clipped squared loss is bounded by $4$ and is $4$-Lipschitz in the
raw prediction.  Symmetrization and contraction bound its expected
uniform deviation by eight times
\eqref{eq:gaussian-rademacher}.  Changing one observation changes that
supremum by at most $4/n$, so bounded-difference concentration gives
\eqref{eq:gaussian-fixed-energy-deviation}.
The suprema are measurable: for each width the relevant loss is
continuous in the parameters, and rational parameter vectors with strict internal-and-readout energy bound are dense in the corresponding closed cylinder.  Taking the
union over the countably many widths gives a countable determining
class.
\end{proof}


\section{Ordinary posterior predictions at a fixed Gaussian scale}
\label{sec:fixed-raw-oracle}

The finite-scale posterior also admits a coefficient-one oracle bound
for its \emph{raw} predictions.  This includes the ordinary posterior
mean, without clipping either the likelihood or the predictor.  For
two or more hidden layers this section assumes $\gamma\geq2$; with one
hidden layer it covers every $\gamma\geq1$.  Sections~\ref{sec:effective-energy} and~\ref{sec:raw-all-scales} treat all
$\gamma\geq1$ at larger depth using a sample-dependent overall prior
scale.

Use the tanh architecture and the full and after-first-layer energies
$Q_N$ and $Q_N^+$ defined above, with $H$
hidden layers, inputs in $[-1,1]^d$, and labels in $[-1,1]$.  In this
section the prior scale is exactly $s^2=N^{-\gamma}$, and the likelihood
variance is $n\lambda s^2$.  Thus
\[
 \Pi(dw)=Z^{-1}\exp[-\widehat R_S(z_w)/(2\lambda s^2)]\,\mu(dw).
\]
A fixed teacher $g$ of width $m$ and weights at most $B\geq1$ has
clones of energy at most
\[
 E=d(d+1)B^2+2Hm^2B^2,\qquad |g|\leq U=mB.
\]
Define
\begin{gather*}
 D_0=d+H+1,\qquad T=8+2H\log2,\qquad
 A_g=\sqrt2mB,\\
 K_g=1+A_g\left(\sqrt2 A_g^{H-1}
                   +\sum_{j=0}^{H-2}A_g^j\right),\\
 G_g(\lambda)=E+1+H\log2
              +2D_0\log\frac{4K_g(U+1)}\lambda .
\end{gather*}
The sum is empty when $H=1$, in which case $K_g=1+2mB$.
We fix a universal constant $C_0\geq1$ for the following smooth-loss
inequality of Srebro, Sridharan, and Tewari~\cite[Theorem~1]{sst}.
For a class $\mathcal F$ on which a nonnegative $s_0$-smooth loss is
bounded by $b$, with probability at least $1-\eta$, simultaneously for
$f\in\mathcal F$,
\begin{align}
 R(f)-\widehat R_S(f)
 &\leq C_0\biggl[\sqrt{\widehat R_S(f)}
   \left(\sqrt{s_0}\log^{3/2}(2n)\mathfrak R_n(\mathcal F)
      +\sqrt{\frac{b\log(1/\eta)}n}\right)\notag\\
 &\hspace{24mm}+s_0\log^3(2n)\mathfrak R_n(\mathcal F)^2
             +\frac{b\log(1/\eta)}n\biggr].
 \label{eq:finite-raw-sst}
\end{align}
Here $\mathfrak R_n$ is worst-case Rademacher complexity.  Increasing
$C_0$ absorbs harmless logarithm conventions.  For raw square loss,
$s_0=2$, and its bound on $\mathcal Z_K$ is $(\sqrt K+1)^2$.
The energy classes are symmetric under readout sign reversal, so the
absolute-value Rademacher convention in that theorem agrees with the
one used here.

For $n\geq2$ and $0<\delta<1$, put
\begin{gather*}
 p=H/2,\qquad J=\lceil\log_2(Tn)\rceil,\qquad
 q_n=\frac1n\log\frac{2(J+1)}\delta,\\
 D_n=\sqrt{2(d+1)\log(n+1)}+\frac{H-1}{2},\qquad
 A_n=\frac{\sqrt2D_n\log^{3/2}(2n)}{H^p\sqrt n},\qquad
 a_n=2^p A_n,\\
 r_n=C_0\left[\sqrt T(A_n+2\sqrt{q_n})+A_n^2+4q_n\right],
 \qquad C_0\geq1,\\
 b_n(g,\delta)=(U+1)\sqrt{\frac{2R(g)\log(2/\delta)}n}
       +\frac{2(U+1)^2\log(2/\delta)}{3n}.
\end{gather*}

\begin{theorem}[Finite-scale raw-output oracle inequality]
\label{thm:fixed-raw-oracle}
Suppose $H\geq2$ and
\[
 \frac1n\leq\lambda\leq1,\qquad
 \lambda^p\geq8C_0a_nT^{p-1/2},\qquad
 \lambda\geq64C_0q_n.
\]
For every fixed teacher $g$, with sample probability at least
$1-\delta$, simultaneously for every $N\geq m$ and $\gamma\geq2$,
\begin{align}
 \int R(z_w)\,\Pi(dw)
 &\leq R(g)+b_n(g,\delta)+\lambda G_g(\lambda)+r_n
       +\frac{16C_0^2Tq_n}{\lambda}
       +2(1+2H)e^{-1/\lambda}.
 \label{eq:fixed-raw-oracle}
\end{align}
The ordinary posterior mean $\int z_w\,d\Pi$ has the same risk bound.
For $H=1$, the statement holds simultaneously for every $\gamma\geq1$
if its calibration and its penultimate term are replaced by
\[
 \frac1n\leq\lambda\leq1,\qquad
 \lambda\geq4C_0(a_n^2+8q_n),\qquad
 \frac{V_n}{\lambda},\quad
 V_n=C_0^2T(a_n+2\sqrt{2q_n})^2.
\]
For a fixed teacher, the resulting rates are
$\widetilde O(n^{-1/H})$ for $H\geq2$ and
$\widetilde O(n^{-1/2})$ for $H=1$, with no learner-width factor.
\end{theorem}

\subsection{The raw posterior tail}

There are $P_N^+=N+(H-1)N^2$ parameters after the first layer, and
the stated regimes give $P_N^+s^2\leq H$.  Let
$F=\widehat R_S(z_w)+\lambda Q_N^+(w)$.
For every sample and every $0<\lambda\leq1$,
\begin{equation}
 \int(1+Q_N^+)\mathbf1_{\{F>T\}}\,d\Pi
     \leq(1+2H)e^{-1/\lambda}.
 \label{eq:fixed-raw-tail}
\end{equation}
Indeed, the centered independent readout gives
$\mathbb E_\mu\widehat R_S(z_w)\leq1+s^2$, so Jensen yields
$Z\geq\exp[-(1+s^2)/(2\lambda s^2)]$.  Split the exponent $F/(2\lambda s^2)$ in half on $F>T$ and discard
the remaining nonnegative empirical loss.  The Gaussian first-layer
factor is left intact.  The resulting Gaussian integral is exact, giving
\begin{align*}
 \int(1+Q_N^+)\mathbf1_{\{F>T\}}\,d\Pi
 &\leq (1+2P_N^+s^2)
 \exp\left[-\frac{T}{4\lambda s^2}
       +\frac{P_N^+}2\log2+\frac{1+s^2}{2\lambda s^2}\right].
\end{align*}
Here $s^2\leq1$, $P_N^+s^2\leq H$, and $\lambda\leq1$ make the
exponent at most $-1/(\lambda s^2)\leq-1/\lambda$ for the stated $T$.
Also $R(z_w)\leq2Q_N^++2$, so the raw risk beyond this objective threshold
is at most twice~\eqref{eq:fixed-raw-tail}.

\subsection{Statistical absorption and the Gaussian comparison}

The smooth-loss bound of Srebro, Sridharan, and Tewari, in the form
stated above, applied to the dyadic energy classes
$\mathcal Z_{2^j}$ for $0\leq j\leq J$ gives a common sample event of
probability at least $1-\delta/2$.  On $F\leq T$, write $k=Q_N^+(w)$;
then $k\leq T/\lambda\leq Tn$ and $\widehat R_S(z_w)\leq T$.
For $k\leq1$ the unit-energy event bounds the deviation by $r_n$.
For $k>1$, the shell of radius at most $2k$ gives
\begin{equation}
 R(z_w)-\widehat R_S(z_w)
 \leq C_0\left[\sqrt T a_nk^p+a_n^2k^{2p}
                   +2\sqrt{2T}\sqrt{kq_n}+8kq_n\right].
 \label{eq:fixed-raw-shell}
\end{equation}
For $p\geq1$, the four terms are at most
\[
 \frac{\lambda k}8,\quad
 \frac{\lambda k}{64C_0},\quad
 \frac{\lambda k}8+\frac{16C_0^2Tq_n}{\lambda},\quad
 \frac{\lambda k}8,
\]
respectively.  For $p=1/2$, the right side of
\eqref{eq:fixed-raw-shell} is
\[
 C_0\sqrt T(a_n+2\sqrt{2q_n})\sqrt k
             +C_0(a_n^2+8q_n)k
 \leq\frac\lambda2 k+\frac{V_n}\lambda.
\]
Thus the respective residual, together with $r_n$, bounds the deviation
above $\widehat R_S+\lambda Q_N^+/2$ throughout $F\leq T$.

The Gaussian entropy inequality gives
\[
 \mathrm{KL}(\nu\Vert\mu)
   \geq\frac1{4s^2}\int Q_N^+\,d\nu-\frac{P_N^+}2\log2.
\]
Posterior variational optimality therefore implies
\begin{align}
 \int[\widehat R_S(z_w)+\lambda Q_N^+/2]\,d\Pi
 &\leq\int\widehat R_S(z_w)\,d\nu
       +2\lambda s^2\mathrm{KL}(\nu\Vert\mu)
       +\lambda H\log2.
 \label{eq:fixed-raw-variation}
\end{align}
Choose $\nu=\mathcal N(w_*,r^2\Sigma)$ around an exact teacher clone,
where $\Sigma$ is the prior covariance,
$\rho=\lambda/[4(U+1)]$, and $r=\rho/K_g$.  The clone's internal
operator norms and its normalized readout norm are at most $A_g$.
For the normalized hidden-layer root mean square errors $e_\ell$,
independence of each perturbation matrix from the preceding layer gives
\[
 e_1\leq\sqrt2r,\qquad
 e_\ell\leq A_ge_{\ell-1}+r,
 \qquad
 (\mathbb E_\nu|z_w(x)-g(x)|^2)^{1/2}\leq rK_g=\rho.
\]
Only $s\sqrt N\leq1$ and $s\leq1$ are used here.  Minkowski gives
$\int\widehat R_S(z_w)d\nu\leq\widehat R_S(g)+\lambda$.
The full parameter count is $P_N=N(d+2)+(H-1)N^2$, with
$P_Ns^2\leq D_0$ in the stated regimes, so the Gaussian entropy formula
gives
\[
 2\lambda s^2\mathrm{KL}(\nu\Vert\mu)
   \leq\lambda E+2\lambda D_0\log(1/r).
\]
Substitution into~\eqref{eq:fixed-raw-variation} bounds its left side
by $\widehat R_S(g)+\lambda G_g(\lambda)$.

Integrate the low-objective deviation, add
\eqref{eq:fixed-raw-tail}, and use nonnegativity to extend the integral
of $\widehat R_S+\lambda Q_N^+/2$ to the whole parameter space.
Bernstein's inequality for the fixed teacher gives
$\widehat R_S(g)\leq R(g)+b_n(g,\delta)$ except on a set of sample
probability $\delta/2$.  This proves the theorem.  All sample events are
independent of $N$ and $\gamma$.  The posterior has bounded density
relative to the Gaussian prior and hence finite second moments of its
predictions; Jensen proves the ordinary-mean assertion.

\subsection{Nearly square-root competitor width in a shallow network}

The shallow theorem can be calibrated to the competitor class instead
of a single fixed teacher.  Fix caps $m_0,B_0\geq1$ before observing the
sample and put
\begin{gather*}
 E_0=d(d+1)B_0^2+2m_0^2B_0^2,\quad U_0=m_0B_0,\quad K_0=1+2m_0B_0,\\
 G_0=E_0+1+H\log2+2D_0\log[4K_0(U_0+1)n],\\
 v_n=\max\{1/n,\,4C_0(a_n^2+8q_n)\},\qquad
 \lambda_n=\max\{\sqrt{V_n/G_0},v_n\}.
\end{gather*}
Whenever $\lambda_n\leq1$, every teacher in the capped class satisfies
\[
 \lambda_nG_g(\lambda_n)+\frac{V_n}{\lambda_n}
       \leq2\sqrt{V_nG_0}+v_nG_0.
\]
The population-risk minimizer over this compact teacher class exists
and has risk at most one, because the zero network is included.  It is
fixed before sampling and can be used in the preceding theorem.  At
fixed $d$ and $\delta$, this proves the coefficient-one oracle bound
\[
 R\left(\int z_w\,d\Pi\right)
 \leq\inf_{\substack{\mathrm{width}(g)\leq m_0\\
                     \mathrm{weights}(g)\leq B_0}}R(g)
 +\widetilde O\left(\frac{m_0B_0+1}{\sqrt n}
                         +\frac{m_0^2B_0^2+1}{n}\right),
\]
with the same statement for posterior-averaged risk.  The notation
$\mathrm{weights}(g)\leq B_0$ means that every absolute weight is at
most $B_0$; the constants may depend on $d$.  The guarantee holds on
one sample event for every $N\geq m_0$ and every $\gamma\geq1$.
Consequently, fixed weight cap and width $m_0\leq n^\alpha$ are allowed
for every $\alpha<1/2$.  If both width and weight cap are at most
$n^\alpha$, every $\alpha<1/4$ is allowed.  Smaller teachers can be
zero-padded to width $m_0$; this is a comparator construction and does
not put a prior on network sizes.

\section{A finite-scale effective energy with no quadratic width cost}
\label{sec:effective-energy}

Write $H\geq2$ for the number of hidden layers and fix
\[
  \sigma^2=t^2N^{-\gamma},\qquad 0<t\leq1,\qquad\gamma\geq1.
\]
The independent Gaussian weight variances are $N\sigma^2/d$ in the
first layer, $\sigma^2$ in each internal layer, and $\sigma^2/N$ in
the readout.  The likelihood variance is $n\lambda\sigma^2$.  Thus the
posterior is the probability law
\begin{equation}
 \Pi(dw)=Z^{-1}\exp\left[-\frac{\widehat R_S(z_w)}{2\lambda\sigma^2}\right]
 \mu(dw).
 \label{eq:effective-posterior}
\end{equation}
The construction below applies at the displayed finite parameters.
No limit in width or prior strength is used.

The ordinary quadratic energy under this prior has an internal-layer
contribution of order $t^2N^{2-\gamma}$.  This is unbounded in width
when $1\leq\gamma<2$, despite the much smaller operator norm of the
Gaussian matrices.  We retain only their singular values above the
Gaussian operator-norm scale when measuring complexity.  This changes
an auxiliary quantity in the proof; it does not restrict or modify the
learner's prior.

\subsection{Spectral truncation and an exponential moment}

Set $a_N=\sigma\sqrt N=tN^{(1-\gamma)/2}$ and $b_N=8a_N$.
The quantity $a_N$ measures the scale of the Gaussian internal
operator norms.  For each internal
matrix $W^\ell$, let $T^\ell=(W^\ell)_{>b_N}$ retain singular values
strictly greater than $b_N$ and set the others to zero.  Let
$\widetilde w$ be the network with these internal matrices and the
original first layer and readout.  Define
\begin{equation}
 q_N(w):=Q_N^+(\widetilde w)
 =\sum_{\ell=2}^H\|T^\ell\|_F^2+N\|v\|_2^2.
 \label{eq:effective-energy}
\end{equation}
This is a nonnegative Borel function and is at most $Q_N^+(w)$.
The first-layer weights are unrestricted in the capacity estimate
for $Q_N^+$ and are unchanged by truncation, so they need not enter
this auxiliary energy.  They still contribute to the prior and to
the teacher comparison below.

\begin{lemma}[Dimension-free cost of the internal Gaussian bulk]
\label{lem:effective-mgf}
Put
\begin{align}
 C_N&=\frac N2\log\frac87+(H-1)\log\frac32,
 \notag\\
 c_H&=\frac12\log\frac87+(H-1)\log\frac32.
 \label{eq:effective-mgf-constant}
\end{align}
Then
\begin{equation}
 \log\mathbb E_\mu\exp\left(\frac{q_N}{16\sigma^2}\right)
 \leq C_N,
 \qquad \sigma^2 C_N\leq a_N^2c_H.
 \label{eq:effective-mgf}
\end{equation}
Consequently, for every probability law $\nu\ll\mu$ of finite
relative entropy,
\begin{equation}
 \operatorname{KL}(\nu\Vert\mu)
 \geq\frac1{16\sigma^2}\int q_N\,d\nu-C_N.
 \label{eq:effective-entropy}
\end{equation}
\end{lemma}
\begin{proof}
The readout contribution, divided by $\sigma^2$,
has a chi-square distribution with $N$ degrees of freedom.
Its exponential moment at parameter $1/16$ is therefore
$(8/7)^{N/2}$.

For an internal matrix $W=\sigma Z$, where $Z$ has independent standard
Gaussian entries, the Gaussian operator-norm bound
\[
 \mathbb P(\|Z\|_{\mathrm{op}}>2\sqrt N+u)\leq e^{-u^2/2}
\]
implies that the operator-norm ball
$B=\{A:\|A\|_{\mathrm{op}}\leq b_N/2\}$ has probability at least
$1/2$.  For example, $u=2\sqrt N$ gives a failure probability at
most $e^{-2N}<1/2$.  The squared Frobenius distance to this ball is
\[
 \operatorname{dist}_F(W,B)^2
   =\sum_j(s_j(W)-b_N/2)_+^2.
\]
Since $s\leq2(s-b_N/2)$ whenever $s>b_N$,
\[
 \|W_{>b_N}\|_F^2\leq4\operatorname{dist}_F(W,B)^2.
\]
Gaussian isoperimetry and $1-\Phi(u)\leq\tfrac12e^{-u^2/2}$ yield
\begin{equation}
 \mathbb P(\|W_{>b_N}\|_F^2>a)
 \leq\frac12\exp\left(-\frac a{8\sigma^2}\right),
 \qquad a>0.
 \label{eq:effective-internal-tail}
\end{equation}
Tail integration now gives
\begin{align*}
 \mathbb E\exp\left(\frac{\|W_{>b_N}\|_F^2}{16\sigma^2}\right)
 &=1+\frac1{16\sigma^2}\int_0^\infty
      e^{a/(16\sigma^2)}
      \mathbb P(\|W_{>b_N}\|_F^2>a)\,da
 \leq\frac32.
\end{align*}
Independence of the layer blocks proves the first inequality
in~\eqref{eq:effective-mgf}.  Its second inequality follows from $\sigma^2N=a_N^2$
and $\sigma^2\leq a_N^2$.
The entropy variational inequality proves
\eqref{eq:effective-entropy}; truncation of $q_N$ justifies the same
conclusion if its integrability has not yet been established.
\end{proof}

\begin{lemma}[Approximation by the truncated network]
\label{lem:effective-approximation}
For every parameter vector and every input,
\begin{align}
 |z_w(x)-z_{\widetilde w}(x)|
 &\leq b_N\sum_{j=1}^{H-1}q_N(w)^{j/2}
 \notag\\
 &\leq8a_N(H-1)(1+q_N(w))^{(H-1)/2}.
 \label{eq:effective-approximation}
\end{align}
Both outputs have absolute value at most $\sqrt{q_N(w)}$.
\end{lemma}
\begin{proof}
The first hidden layers agree.  Write
$e_\ell=\|h^\ell-\widetilde h^\ell\|_2/\sqrt N$.
Using
\[
 W^\ell h^{\ell-1}-T^\ell\widetilde h^{\ell-1}
 =T^\ell(h^{\ell-1}-\widetilde h^{\ell-1})
      +(W^\ell-T^\ell)h^{\ell-1},
\]
the contraction of tanh, $\|h^{\ell-1}\|_2\leq\sqrt N$, and
$\|T^\ell\|_{\mathrm{op}}\leq\sqrt{q_N}$ give
\[
 e_1=0,\qquad e_\ell\leq\sqrt{q_N}\,e_{\ell-1}+b_N.
\]
Multiplication by $\sqrt N\|v\|_2\leq\sqrt{q_N}$ gives the first
bound.  The second follows from $b_N=8a_N$ and
$q_N^{j/2}\leq(1+q_N)^{(H-1)/2}$.  Bounded hidden activations give
the output bounds.  In particular, no operator-norm contribution
$b_N$ is needed in the multiplying gain of the recursion.
\end{proof}

The posterior variational identity and
Lemma~\ref{lem:effective-mgf} imply, for every comparison law $\nu$
of finite relative entropy,
\begin{align}
 \int\left[\widehat R_S(z_w)+\frac\lambda8q_N(w)\right]\Pi(dw)
 &\leq\int\widehat R_S(z_w)\,\nu(dw)
       +2\lambda\sigma^2\operatorname{KL}(\nu\Vert\mu)
       +2\lambda\sigma^2 C_N.
 \label{eq:effective-variational}
\end{align}
All these quantities are finite for the comparison law constructed
next and for the posterior.  For the latter, its density relative to
$\mu$ is bounded by $1/Z$, and $q_N\leq Q_N$ has all Gaussian moments.

\subsection{A teacher neighborhood with linear width cost}

Let $g$ have width $m$, depth $H+1$, and weights bounded by $B\geq1$.
At any $N\geq m$, balanced cloning computes $g$ exactly and has energy
at most
\begin{equation}
 E=d(d+1)B^2+2Hm^2B^2,
 \qquad |g(x)|\leq U:=mB.
 \label{eq:effective-teacher-energy}
\end{equation}
The neighborhood below narrows only the noise components that act on
the teacher's group-constant activations.  The other Gaussian
coordinates remain present.

\begin{lemma}[Teacher neighborhood at every finite width and scale]
\label{lem:effective-teacher-tube}
For $0<\rho\leq1$, define
\begin{align*}
 b_0&=2\sqrt{2[2\log9+\log(4(H+1))]},
 & A_0&=2+3mB+b_0,\\
 P_0&=1+2mB(d+H+2)A_0^{H-1},
 & \eta&=\min\{1/2,\rho/(2P_0)\},\\
 D_g&=d+2+(H-1)m,
 & c_0&=3/8.
\end{align*}
Let $U_\rho(g)$ be the event that $|z_w(x)-g(x)|<\rho$ at every
$x\in[-1,1]^d$.  Then, for $N\geq m$, $\gamma\geq1$, and $0<t\leq1$,
\begin{equation}
 -\log\mu(U_\rho(g))
 \leq\frac E{2\sigma^2}+ND_g\log\frac1{c_0\eta}+\log2.
 \label{eq:effective-teacher-tube}
\end{equation}
In particular, with $\rho=\lambda/[4(U+1)]$ and $0<\lambda\leq1$,
there is a comparison law $\nu$ with
\begin{align}
 \int\widehat R_S(z_w)\,d\nu
    +2\lambda\sigma^2\operatorname{KL}(\nu\Vert\mu)
 &\leq\widehat R_S(g)+\lambda E+\lambda
 \notag\\
 &\quad+2\lambda a_N^2
       \left[D_g\log\frac1{c_0\eta}+\log2\right].
 \label{eq:effective-comparison}
\end{align}
\end{lemma}
\begin{proof}
Let $w_*$ be the balanced clone and let $P$ project onto its
$m$-dimensional group-constant subspace, with orthonormal basis
$u_j=k_j^{-1/2}\mathbf1_{\text{group }j}$.  In standardized Gaussian
coordinates impose $|\xi|\leq\eta$ on each first-layer and readout
coordinate, and
$|\langle Z_{i,*},u_j\rangle|\leq\eta$ for every internal row and
group.  There are $ND_g$ independent scalar constraints.  A centered
standard Gaussian has probability at least $c_0\eta$ of lying in
$[-\eta,\eta]$ when $0<\eta\leq1$; for example, combine symmetric
unimodality with Chebyshev's probability $3/4$ for $[-2,2]$.
The cylinder therefore has probability at least $(c_0\eta)^{ND_g}$.

Conditioning on this centrally symmetric convex cylinder preserves
centered unit-subgaussian linear forms.  This follows directly by
factorizing into Gaussian coordinates along the orthonormal constrained
subspaces, or from the strongly log-concave moment-generating-function
bound.  Two sphere nets of cardinality at most $9^N$ consequently give
\[
 \mathbb P(\|Z\|_{\mathrm{op}}>2u\sqrt N\mid\text{cylinder})
 \leq2\exp[(2\log9-u^2/2)N].
\]
Intersecting with $\|Z\|_{\mathrm{op}}\leq b_0\sqrt N$ for every
internal layer leaves a symmetric event $A$ with probability at least
$\tfrac12(c_0\eta)^{ND_g}$.

Translate by the standardized clone.  Write
$a_N=\sigma\sqrt N=tN^{(1-\gamma)/2}\leq1$.
The internal perturbation is
$a_NZ/\sqrt N=a_NZP/\sqrt N+a_NZ(I-P)/\sqrt N$.
Its second term annihilates the clone's activations, while the first
has operator norm at most $\eta\sqrt m$.
For the normalized hidden errors this gives, uniformly on the cube,
\[
 e_1\leq\eta(d+1),\qquad
 e_\ell\leq(\sqrt2mB+b_0)e_{\ell-1}+\eta\sqrt m.
\]
Thus $e_H\leq\eta(d+H+2)A_0^{H-1}$.
The clone readout has $\sqrt N\|v_*\|_2\leq\sqrt2mB$, and the
perturbation readout has $\ell^1$ norm at most $\eta$.
The output error is at most $\eta P_0\leq\rho/2$.
Hence the translated event is contained in $U_\rho(g)$.

For standard Gaussian measure $G$ and a centrally symmetric event $A$,
\[
 G(A+a_*)\geq G(A)e^{-\|a_*\|_2^2/2}.
\]
Indeed, the Gaussian density ratio averaged over $A$ is
$e^{-\|a_*\|_2^2/2}\mathbb E[e^{-\langle Z,a_*\rangle}\mid A]$,
and its second factor is at least one by symmetry and Jensen.
Here $\|a_*\|_2^2=Q_N(w_*)/\sigma^2\leq E/\sigma^2$.
This proves~\eqref{eq:effective-teacher-tube}.

For the last assertion, use $\nu=\mu(\cdot\mid U_\rho(g))$.
Its relative entropy is $-\log\mu(U_\rho(g))$, and
\[
 \int\widehat R_S(z_w)\,d\nu
 \leq\widehat R_S(g)+2\rho(U+1)+\rho^2
 \leq\widehat R_S(g)+\lambda.
\]
Finally $\sigma^2N=a_N^2$ and
$\sigma^2\leq a_N^2$ give~\eqref{eq:effective-comparison}.
\end{proof}

\section{Raw predictions throughout the mean-field-or-stronger regime}
\label{sec:raw-all-scales}

We now prove a finite-posterior oracle bound for the ordinary raw-output
posterior mean at every $\gamma\geq1$ and every eligible width, at any
fixed depth.  The overall prior scale decreases explicitly with sample
size.  Thus this theorem does not claim the same result at a fixed
prefactor $t=1$.

Use the model~\eqref{eq:effective-posterior}, with $H\geq2$, and the
effective energy $q_N$ from~\eqref{eq:effective-energy}.  Write
$a_N=\sigma\sqrt N$; the theorem below applies to any finite
$\sigma>0$ satisfying its stated bound on $a_N$.  To apply the preceding
lemmas, represent the same Gaussian law with $\gamma=1$ and $t=a_N$.
Recall the constants
$C_N,c_H$ from
\eqref{eq:effective-mgf-constant}.  Put
\begin{gather}
 T=64(1+c_H),\quad p=H/2,\quad J=\lceil\log_2(Tn)\rceil,
 \quad u_n=\frac1n\log\frac{2(J+1)}\delta,\notag\\
 D_n=\sqrt{2(d+1)\log(n+1)}+\frac{H-1}{2},\quad
 A_n=\frac{\sqrt2D_n\log^{3/2}(2n)}{H^p\sqrt n},\quad a_n=2^pA_n,
 \notag\\
 r_n=C_0[2\sqrt T(A_n+2\sqrt{u_n})+A_n^2+4u_n].
 \label{eq:raw-all-statistics}
\end{gather}
For each eligible $\lambda$, also define
\begin{equation}
 \nu_H=\frac{H-1}{2},\qquad b_H=\max\{1,\nu_H\},\qquad
 M_n=5T+r_n+\frac{512C_0^2Tu_n}{\lambda}.
 \label{eq:raw-all-risk-budget}
\end{equation}
Here $C_0\geq1$ is the same universal smooth-loss constant as in
\eqref{eq:finite-raw-sst}.  For a fixed teacher $g$ use $E,U,D_g,P_0,c_0$
from Lemma~\ref{lem:effective-teacher-tube}, and set
\[
 \rho=\frac\lambda{4(U+1)},\qquad
 \eta=\min\{1/2,\rho/(2P_0)\}.
\]
The fixed-teacher Bernstein term is
\[
 b_n(g,\delta)=(U+1)\sqrt{\frac{2R(g)\log(2/\delta)}n}
        +\frac{2(U+1)^2\log(2/\delta)}{3n}.
\]

\begin{theorem}[An all-width, all-exponent raw-output oracle bound]
\label{thm:raw-all-scales}
Let $H\geq2$, $n\geq2$, and $0<\delta<1$.  Suppose
\begin{gather}
 \frac1n\leq\lambda\leq1,\qquad
 \lambda^p\geq128C_0a_nT^{p-1/2},\qquad
 \lambda\geq512C_0u_n,
 \label{eq:raw-all-lambda}\\
 0<a_N\leq\min\left\{1,
       \frac{\lambda^{b_H}}
       {768\sqrt{M_n}(H-1)2^{\nu_H}T^{\max\{\nu_H-1,0\}}}\right\}.
 \label{eq:raw-all-prefactor}
\end{gather}
For every fixed width-$m$, weight-$B$ teacher $g$, with probability at
least $1-\delta$ over the sample, simultaneously for every $N\geq m$,
and every $\sigma>0$ satisfying~\eqref{eq:raw-all-prefactor},
\begin{align}
 \int R(z_w)\,\Pi(dw)
 &\leq R(g)+b_n(g,\delta)
  +\lambda\left[E+1+2a_N^2\left\{c_H
                 +D_g\log\frac1{c_0\eta}+\log2\right\}\right]
       \notag\\
 &\quad+r_n+\frac{512C_0^2Tu_n}{\lambda}
              +48\sqrt{M_n}a_N(H-1)2^{\nu_H}+66e^{-1/\lambda}.
 \label{eq:raw-all-oracle}
\end{align}
The same right side bounds the risk of the ordinary posterior mean.
It also bounds the usual clipped prediction conventions.
No limit in width or prior scale is taken.
\end{theorem}

For sufficiently large $n$, an explicit calibration is
\begin{gather}
 \lambda_n=\max\left\{1/n,
            [128C_0a_nT^{p-1/2}]^{1/p},512C_0u_n\right\},\qquad
 t_n=\frac{\lambda_n^{b_H}}
  {768\sqrt{M_n}(H-1)2^{\nu_H}T^{\max\{\nu_H-1,0\}}}.
 \label{eq:raw-all-schedule}
\end{gather}
For every $\gamma\geq1$ and every $N\geq m$, the prior
$\sigma^2=t_n^2N^{-\gamma}$ satisfies~\eqref{eq:raw-all-prefactor},
because $a_N=t_nN^{(1-\gamma)/2}\leq t_n$.  The same $t_n$ therefore
works simultaneously at all these widths and exponents.
At fixed $d,H,\delta$ these have orders
$\widetilde O(n^{-1/H})$ and $\widetilde O(n^{-b_H/H})$,
respectively; $M_n$ remains bounded as $n$ grows with these resources
fixed.  Thus the sufficient prefactor has exponent $1/2$ at $H=2$,
$1/3$ at $H=3$, and $(H-1)/(2H)$ for $H\geq4$.
The excess in~\eqref{eq:raw-all-oracle} is
$\widetilde O(n^{-1/H})$ for a fixed teacher.  The calibration depends
on $n,d,H,\delta$, not on the unknown data law or on a selected teacher.

\subsection{A weighted posterior tail for effective energy}

Let $F(w)=\widehat R_S(z_w)+\lambda q_N(w)$.  Every sample satisfies
\begin{equation}
 \int(1+q_N)\mathbf1_{\{F>T\}}\,d\Pi
       \leq33e^{-1/(\lambda\sigma^2)}\leq33e^{-1/\lambda}.
 \label{eq:raw-effective-tail}
\end{equation}
To prove this, the centered readout gives
$\mathbb E_\mu\widehat R_S\leq1+\sigma^2$ and hence
$Z\geq\exp[-(1+\sigma^2)/(2\lambda\sigma^2)]$.
Write $\beta=1/(2\lambda\sigma^2)$.  On $F>T$,
\[
 e^{-\beta\widehat R_S}
   \leq e^{-\beta T/16}e^{q_N/(32\sigma^2)},\qquad
 (1+q_N)e^{q_N/(32\sigma^2)}
    \leq(1+32\sigma^2)e^{q_N/(16\sigma^2)}
    \leq33e^{q_N/(16\sigma^2)}.
\]
The moment bound~\eqref{eq:effective-mgf} therefore makes the left
side of~\eqref{eq:raw-effective-tail} at most
\[
 33\exp\left(C_N+\frac{1+\sigma^2}{2\lambda\sigma^2}
                         -\frac{T}{32\lambda\sigma^2}\right).
\]
Since $\sigma^2C_N\leq c_H$, $\sigma^2\leq1$, and $\lambda\leq1$,
the exponent is at most
\[
 \frac{2c_H+2-T/16}{2\lambda\sigma^2}
       \leq-\frac1{\lambda\sigma^2}.
\]
This proves the claim.  Hard spectral truncation leaves the readout
unchanged, so both original and truncated outputs have absolute value
at most $\sqrt{q_N}$.  Thus the raw population risk on $F>T$ contributes
at most $66e^{-1/\lambda}$.

\subsection{Smooth-loss absorption on the low-objective region}

Apply~\eqref{eq:finite-raw-sst} to the deterministic energy classes
$\mathcal Z_{2^j}$, $0\leq j\leq J$, with failure probabilities
$\delta/[2(J+1)]$.  The common event has probability at least
$1-\delta/2$ and does not depend on $N$ or $\sigma$.
On $F\leq T$ write $q=q_N(w)$ and let $\widetilde w$ denote the
truncated network.  Then $q\leq T/\lambda\leq Tn$,
$\widehat R_S(z_w)\leq T$, and
Lemma~\ref{lem:effective-approximation} gives
\[
 e_w:=\|z_w-z_{\widetilde w}\|_\infty
       \leq8a_N(H-1)(1+q)^{\nu_H}\leq\frac{T}{96\sqrt{M_n}}
       \leq\sqrt T.
\]
For this use~\eqref{eq:raw-all-prefactor},
$1+T/\lambda\leq2T/\lambda$, $\lambda\leq1$, and $M_n\geq T$.
In particular, Minkowski yields
$\sqrt{\widehat R_S(z_{\widetilde w})}\leq2\sqrt T$.

For $q\leq1$, the unit-energy smooth-loss bound is at most $r_n$.
For $q>1$, a dyadic shell of radius at most $2q$ gives
\begin{align*}
 R(z_{\widetilde w})-\widehat R_S(z_{\widetilde w})
 \leq C_0[2\sqrt T a_nq^p+a_n^2q^{2p}
                  +4\sqrt{2T}\sqrt{qu_n}+8qu_n].
\end{align*}
The four terms are at most
\[
 \frac{\lambda q}{64},\quad
 \frac{\lambda q}{16384C_0},\quad
 \frac{\lambda q}{64}+\frac{512C_0^2Tu_n}{\lambda},\quad
 \frac{\lambda q}{64},
\]
respectively, by~\eqref{eq:raw-all-lambda}, $p\geq1$, and
$q\leq T/\lambda$.  Their sum is at most
$\lambda q/16+512C_0^2Tu_n/\lambda$.

For both energy cases, the preceding inequality implies
\[
 R(z_{\widetilde w})\leq
 \widehat R_S(z_{\widetilde w})+\lambda q/16
             +r_n+512C_0^2Tu_n/\lambda\leq M_n.
\]
Here $\widehat R_S(z_{\widetilde w})\leq4T$ and $\lambda q\leq T$.
Minkowski therefore transfers the two raw losses with cost at most
\begin{align*}
 R(z_w)-R(z_{\widetilde w})
  +\widehat R_S(z_{\widetilde w})-\widehat R_S(z_w)
 &\leq2e_w(\sqrt{M_n}+\sqrt T)+2e_w^2\\
 &\leq6\sqrt{M_n}e_w\\
 &\leq48\sqrt{M_n}a_N(H-1)(1+q)^{\nu_H}.
\end{align*}
If $q>1$ and $\nu_H\geq1$, use
$(1+q)^{\nu_H}\leq2^{\nu_H}(T/\lambda)^{\nu_H-1}q$;
if $0<\nu_H<1$, use $(1+q)^{\nu_H}\leq2^{\nu_H}q$.
In either case~\eqref{eq:raw-all-prefactor} bounds the transfer cost
by $\lambda q/16$.  For $q\leq1$ its bound is
$48\sqrt{M_n}a_N(H-1)2^{\nu_H}$.
Consequently, throughout $F\leq T$,
\begin{equation}
 R(z_w)\leq\widehat R_S(z_w)+\frac\lambda8 q_N(w)
       +r_n+\frac{512C_0^2Tu_n}{\lambda}
       +48\sqrt{M_n}a_N(H-1)2^{\nu_H}.
 \label{eq:raw-effective-deviation}
\end{equation}

\begin{proof}[Proof of Theorem~\ref{thm:raw-all-scales}]
Integrate~\eqref{eq:raw-effective-deviation} over $F\leq T$ and use
\eqref{eq:raw-effective-tail} outside.  The integrand
$\widehat R_S+\lambda q_N/8$ is nonnegative, so its restricted integral
is at most its full integral.  Apply the variational comparison
\eqref{eq:effective-variational} with the teacher-cylinder comparison
from~\eqref{eq:effective-comparison}.  The moment constant satisfies
$\sigma^2C_N\leq a_N^2c_H$.  These statements give the right side of
\eqref{eq:raw-all-oracle} with $\widehat R_S(g)$ in place of
$R(g)+b_n(g,\delta)$.

Bernstein's inequality for the fixed teacher supplies that replacement
except on a sample set of probability $\delta/2$.  The resulting event
is common to all the claimed widths and scales.  The posterior has
bounded density relative to its Gaussian prior, and
$|z_w|^2\leq Q_N(w)$; thus the ordinary posterior mean is integrable,
and Jensen proves its risk bound.  Clipping cannot increase squared
loss for labels in $[-1,1]$, giving the remaining conventions.
\end{proof}

The general scale condition also covers a fixed prefactor $t=1$
when $\gamma>1$ and
\[
 N\geq\left(\frac{768\sqrt{M_n}(H-1)2^{\nu_H}
                          T^{\max\{\nu_H-1,0\}}}
                         {\lambda_n^{b_H}}\right)^{2/(\gamma-1)}.
\]
At fixed $d,H,\delta,\gamma$, this sufficient width is
$\widetilde O(n^{2b_H/[H(\gamma-1)]})$.  It is an additional large-width
corollary, not an all-width fixed-prefactor theorem for
$1\leq\gamma<2$.

The coefficient of $R(g)$ is one.  A teacher may therefore be a
population-risk minimizer over any specified compact bounded-weight
class, chosen before sampling.  At fixed $d,H,\delta$, the theorem
competes with every such class whose width and weight caps are at most
$n^\alpha$, for every $\alpha<1/(4H)$: its excess tends to zero,
uniformly over every learner width above the class width and every
$\gamma\geq1$.  With a fixed weight cap the allowed range is
$\alpha<1/(2H)$.  This follows from $E=O(m^2B^2+B^2)$,
$U=mB$, the polynomial teacher-cylinder constants, and the fact that
the population minimizer has risk at most one.

\subsection{The first layer need not become more concentrated}
\label{sec:first-layer-scale-extension}

The finite-scale theorem also permits a conventional fixed fan-in
variance in the first layer.  This is useful because the effective
energy and its capacity bound leave the first-layer weights unrestricted.

\begin{corollary}[A fixed first-layer Gaussian scale]
\label{cor:first-layer-scale-extension}
In Theorem~\ref{thm:raw-all-scales}, replace the first-layer variance
$N\sigma^2/d=a_N^2/d$ by $v_1^2/d$, where $a_N\leq v_1\leq1$.
Leave the internal and readout priors and the likelihood unchanged.
Then the same oracle bound holds, on the same sample event, with the
same displayed constants.  In particular, one may use variance $1/d$
in the first layer for all sample sizes and widths.
\end{corollary}
\begin{proof}
The effective energy $q_N$ uses only the internal matrices and readout,
so its Gaussian exponential moment is unchanged.  The statistical event
allows arbitrary first-layer weights, and the first layer is never
modified by spectral truncation.  The prior evidence lower bound
also remains unchanged, because it uses only the centered independent
readout and bounded hidden activations.

For the teacher cylinder, first-layer noise now has normalized amplitude
$v_1\leq1$, so the bound $e_1\leq\eta(d+1)$ remains valid.
The normalized translation cost of the first-layer clone becomes
\[
 \frac{d\sigma^2}{v_1^2}\|W_*^1\|_F^2
 \leq d(d+1)B^2\frac{a_N^2}{v_1^2}
 \leq d(d+1)B^2.
\]
All other translation costs and cylinder probabilities are unchanged.
Thus the same teacher comparison applies.  The proof of
Theorem~\ref{thm:raw-all-scales} uses only these inputs.
\end{proof}

\section{A raw-output PAC--Bayes bound for smaller widths}
\label{sec:raw-small-width}

A second finite-posterior argument complements spectral compression.
It is useful when the Gaussian inverse scale is modest relative to the
sample count.  Unlike a direct bounded-loss Hoeffding argument, it uses
the posterior's own empirical objective to avoid one inverse power of
$\lambda$.

Fix a learner width $N$ and a Gaussian scale $\sigma>0$ with
$N\sigma^2\leq1$.  The likelihood variance remains
$n\lambda\sigma^2$, where $0<\lambda\leq1$.  For a teacher $g$ use the
constants in Lemma~\ref{lem:effective-teacher-tube} and put
\begin{equation}
 \rho=\frac\lambda{4(U+1)},\quad
 \eta=\min\{1/2,\rho/(2P_0)\},\quad
 \mathcal A_g(\lambda)=\frac E2
          +D_g\log\frac1{c_0\eta}+\log2.
 \label{eq:raw-small-Ag}
\end{equation}
The comparison~\eqref{eq:effective-comparison} implies that the
posterior's raw empirical objective is at most
$\widehat R_S(g)+\lambda+2\lambda\mathcal A_g(\lambda)$.

\begin{lemma}[Raw-output PAC--Bayes comparison]
\label{lem:raw-small-width}
Let $N\sigma^2\leq1$, $0<\lambda\leq1$, and $0<\delta'<1$.  Suppose
\begin{equation}
 \frac{256}{n\lambda^2\sigma^2}\leq1.
 \label{eq:raw-small-calibration}
\end{equation}
With sample probability at least $1-\delta'$,
\begin{align}
 \int R(z_w)\,\Pi(dw)
 &\leq\widehat R_S(g)+\lambda+2\lambda\mathcal A_g(\lambda)
       +4\lambda\sigma^2\log(1/\delta')\notag\\
 &\quad+\frac{512}{n\lambda^2\sigma^2}
                   +\sqrt{24}e^{-7/(\lambda\sigma^2)}.
 \label{eq:raw-small-width-bound}
\end{align}
The same bound holds for the ordinary posterior mean.
\end{lemma}

\begin{proof}
Let $M^2=128/\lambda$, and write
$z_w^{[M]}=\max\{-M,\min\{M,z_w\}\}$.  Its square loss is bounded by
$b=(M+1)^2\leq512/\lambda$.  Write $\beta=1/(2\lambda\sigma^2)$ and
$u=\beta b/n\leq1$.

For any random loss $\ell\in[0,b]$, convexity of the exponential gives
\[
 \mathbb E e^{-u\ell/b}
 \leq1-(1-e^{-u})\mathbb E\ell/b
 \leq\exp[-(1-e^{-u})\mathbb E\ell/b].
\]
Apply this independently to the $n$ observations, integrate over the
prior, and use Markov's inequality and the entropy variational
inequality.  On an event of probability at least $1-\delta'$, every
comparison law $\nu$ satisfies
\[
 \int R(z_w^{[M]})\,d\nu
 \leq\frac{u}{1-e^{-u}}\left[
       \int\widehat R_S(z_w^{[M]})\,d\nu
       +\frac{\mathrm{KL}(\nu\Vert\mu)+\log(1/\delta')}{\beta}
                          \right].
\]
Since $u/(1-e^{-u})\leq1+u$, put
\[
 A_\Pi=\int\widehat R_S(z_w^{[M]})\,d\Pi
                  +\frac{\mathrm{KL}(\Pi\Vert\mu)}\beta.
\]
Clipping to $[-M,M]$ cannot increase squared loss because $M\geq1$.
Posterior variational optimality, compared with the prior itself,
therefore gives
\[
 0\leq A_\Pi
 \leq\int\widehat R_S(z_w)\,d\Pi
                  +\frac{\mathrm{KL}(\Pi\Vert\mu)}\beta
 \leq\int\widehat R_S(z_w)\,d\mu
 \leq1+\sigma^2\leq2.
\]
Consequently
\begin{align*}
 \int R(z_w^{[M]})\,d\Pi
 &\leq A_\Pi+2u+\frac{2\log(1/\delta')}{\beta}\\
 &\leq\widehat R_S(g)+\lambda+2\lambda\mathcal A_g(\lambda)
       +\frac{512}{n\lambda^2\sigma^2}
       +4\lambda\sigma^2\log(1/\delta').
\end{align*}
This step does not assume a low-noise or Bernstein condition on the
unknown data law.

It remains to remove the auxiliary clipping.  Set $r=N\|v\|_2^2$.
Under the prior, $r/\sigma^2$ has the chi-square law with $N$ degrees
of freedom.  Thus $\mathbb E_\mu r\leq1$,
$\mathbb E_\mu r^2\leq3$, and
$R(z_w)\leq2r+2$ gives $\mathbb E_\mu R(z_w)^2\leq24$.
The elementary exponential-moment bound for a chi-square variable
implies
\[
 \mu\{r>M^2\}\leq e^{-M^2/(8\sigma^2)},
\]
since $M^2/\sigma^2\geq3N$.  The centered readout and Jensen give
$Z\geq e^{-1/(\lambda\sigma^2)}$.  As the likelihood is at most one,
Cauchy--Schwarz yields
\begin{align*}
 \int R(z_w)\mathbf1_{\{r>M^2\}}\,d\Pi
 &\leq\sqrt{24}\exp\left[\frac1{\lambda\sigma^2}
                              -\frac{M^2}{16\sigma^2}\right]
 =\sqrt{24}e^{-7/(\lambda\sigma^2)}.
\end{align*}
On $r\leq M^2$ the raw and auxiliary predictions coincide at every
input.  Adding this tail proves~\eqref{eq:raw-small-width-bound}.
The finite posterior second moments also justify Jensen's inequality
for the ordinary posterior mean.
\end{proof}

For the fixed-prefactor prior $\sigma^2=N^{-\gamma}$, $\gamma\geq1$,
the leading sample term in this lemma is
$512N^\gamma/(n\lambda^2)$.  To obtain one event for all widths at a
fixed $\gamma$, take $\delta'_N=\delta/(4N^2)$ and use a union bound.
The confidence term remains uniform because
\[
 \frac{\log(4N^2/\delta)}{N^\gamma}
       \leq\log(4/\delta)+1.
\]
A separate fixed-teacher concentration event replaces
$\widehat R_S(g)$ by its population risk and Bernstein remainder.

\section{Fixed-prefactor priors at every exponent above one}
\label{sec:fixed-above-one-raw}

The small-width PAC--Bayes bound and the large-width spectral bound
also give a raw-output oracle inequality at every fixed
$\gamma>1$, with the fixed prefactor $t=1$.  Its rate is slower
than Theorem~\ref{thm:raw-all-scales}, but no sample-dependent change
to the prior is needed.  For $\gamma\geq2$, Theorem~\ref{thm:fixed-raw-oracle}
already gives the better rate and should be used instead.

Fix $H\geq2$ and $\gamma>1$, and put
\begin{equation}
 b_H=\max\{1,(H-1)/2\},\qquad
 \kappa_{H,\gamma}=
   \frac{\gamma-1}{3(\gamma-1)+2b_H\gamma},\qquad
 \lambda_n=n^{-\kappa_{H,\gamma}}.
 \label{eq:fixed-above-one-rate}
\end{equation}
The prior scale is $\sigma^2=N^{-\gamma}$ and the likelihood variance
is $n\lambda_nN^{-\gamma}$.

\begin{theorem}[A raw oracle bound for every fixed $\gamma>1$]
\label{thm:fixed-above-one-raw}
Let inputs lie in $[-1,1]^d$, labels lie in $[-1,1]$, and let $g$ be a
fixed width-$m$, weight-$B$ teacher.  For all sufficiently large $n$,
with probability at least $1-\delta$ over the sample, the actual
finite posterior satisfies, simultaneously for every $N\geq m$,
\begin{equation}
 \int R(z_w)\,\Pi_{N,\gamma,S}(dw)
 \leq R(g)+\widetilde O_{d,H,\gamma,\delta,g}
                         (n^{-\kappa_{H,\gamma}}).
 \label{eq:fixed-above-one-oracle}
\end{equation}
The same result holds for the ordinary posterior mean.  More
explicitly, the excess is bounded by
\begin{equation}
 C_{d,H,\gamma,\delta}\,n^{-\kappa_{H,\gamma}}
 \left[E+1+D_g\log\bigl((n+1)P_0(U+1)\bigr)\right]
       +b_n(g,\delta/2),
 \label{eq:fixed-above-one-teacher}
\end{equation}
for a finite constant independent of $m,B,N$ and the data law, after a
sample threshold depending only on $d,H,\gamma,\delta$.
The teacher quantities $E,U,D_g,P_0$ are those of
Lemma~\ref{lem:effective-teacher-tube}.
\end{theorem}

\begin{proof}
Write $\kappa=\kappa_{H,\gamma}$.  Directly from its definition,
$0<\kappa<1/H\leq1/2$.  Hence $\lambda_n$ satisfies
\eqref{eq:raw-all-lambda} for sufficiently large $n$, using
$\delta/2$ in that theorem.  Its risk budget $M_n$ is then bounded
above and below by positive constants depending only on $d,H,\delta$.
The threshold in the fixed-prefactor corollary after
Theorem~\ref{thm:raw-all-scales} is therefore at most
\[
 W_n=C_{d,H,\gamma,\delta}\lambda_n^{-2b_H/(\gamma-1)}
\]
for a sufficiently large constant $C_{d,H,\gamma,\delta}$.  Every integer
width $N\geq\max\{m,W_n\}$ satisfies the scale condition of that
theorem.  Its bound is at most~\eqref{eq:fixed-above-one-teacher},
since the statistical residuals $\widetilde O(n^{-1/2})$ and
$\widetilde O(n^{-1}/\lambda_n)$ are faster than $n^{-\kappa}$,
and the spectral transfer residual is $O(\lambda_n^{b_H})$.

For $m\leq N<W_n$, apply Lemma~\ref{lem:raw-small-width} with
$\sigma^2=N^{-\gamma}$ and $\delta'_N=\delta/(4N^2)$.  Its sample
term is bounded uniformly by
\begin{align*}
 \frac{512W_n^\gamma}{n\lambda_n^2}
 &\leq C_{d,H,\gamma,\delta}\,
 n^{-1+\kappa(2+2b_H\gamma/(\gamma-1))}
 = C_{d,H,\gamma,\delta}\,n^{-\kappa}.
\end{align*}
The last equality is exactly~\eqref{eq:fixed-above-one-rate}.  In
particular the calibration~\eqref{eq:raw-small-calibration} holds
throughout these widths once $n$ is sufficiently large.
The confidence term is uniform because
\[
 N^{-\gamma}\log(4N^2/\delta)\leq\log(4/\delta)+1.
\]
The cylinder cost $\mathcal A_g(\lambda_n)$ is at most a constant
multiple of $E+1+D_g\log((n+1)P_0(U+1))$, so the remaining terms of
the lemma have the asserted form.

The large-width proof uses one sample event of probability at least
$1-\delta/2$, including its fixed-teacher Bernstein event.  Intersect
it with the small-width PAC--Bayes events.  Since
$\sum_{N\geq1}\delta/(4N^2)<\delta/2$, the total probability is at
least $1-\delta$.  The teacher concentration event supplies
$b_n(g,\delta/2)$ at the smaller widths as well.  Every width is
covered, and Jensen gives the ordinary-mean assertion.
\end{proof}

\subsection{Calibrating to a growing competitor class}

The fixed-teacher calibration is not optimal for a growing class.  Put
\[
 p_{H,\gamma}=2+\frac{2b_H\gamma}{\gamma-1},
 \qquad \kappa_{H,\gamma}=\frac1{p_{H,\gamma}+1}.
\]
The proof above gives, more generally, a leading tradeoff
\[
 \widetilde O\left(\lambda E_0+
                  \frac1{n\lambda^{p_{H,\gamma}}}\right)
\]
for a specified competitor-energy cap $E_0$, whenever
$\lambda=n^{-\theta}$ with $0<\theta<1/p_{H,\gamma}$.
Indeed $p_{H,\gamma}>H$, so this range also satisfies the spectral
statistical calibration $\theta<1/H$.  Other terms in the proof are
smaller at the calibrations below.

For width and weight caps both at most $n^\alpha$, choose
\[
 \alpha<\frac1{4p_{H,\gamma}},\qquad
 \theta=\frac{1+4\alpha}{p_{H,\gamma}+1},\qquad
 \lambda_n=n^{-\theta}.
\]
Then $4\alpha<\theta<1/p_{H,\gamma}$.  A population-risk minimizer in
the compact capped class has risk at most one, since zero is included.
The same all-width proof yields the coefficient-one oracle excess
\[
 \widetilde O\left(n^{-(1-4p_{H,\gamma}\alpha)/(p_{H,\gamma}+1)}\right).
\]
The fixed-teacher Bernstein term is faster, and the cylinder logarithm
is smaller than the energy term up to logarithms.  With a fixed weight
cap and only width growing as $n^\alpha$, the corresponding range and
rate are
\[
 \alpha<\frac1{2p_{H,\gamma}},\qquad
 \lambda_n=n^{-(1+2\alpha)/(p_{H,\gamma}+1)},\qquad
 \widetilde O\left(n^{-(1-2p_{H,\gamma}\alpha)/(p_{H,\gamma}+1)}\right).
\]
These calibrations use the specified class caps, not the unknown data
law or a selected minimizer.  The prior remains a single dense
Gaussian at each width; no mixture over architectures is used.

\section{A Gaussian mean-field cover from entropy balls}
\label{sec:mf-entropy-cover}

This section concerns the unmodified mean-field Gaussian prior:
first-layer variance $1/d$, internal variance $1/N$, and readout
variance $1/N^2$.  The learner has $H\geq2$ hidden tanh layers.
We construct a finite cover of most prior outputs.  The construction
uses an auxiliary Gaussian measure on one deterministic finite
architecture solely to count separated functions.  This auxiliary
measure is not the learner's prior and introduces no mixture over
learner architectures.

For a finite set of $s$ inputs and a positive semidefinite covariance
matrix $Q$, let $P_Q$ be the law of $h=\tanh Z$ for
$Z\sim\mathcal N(0,Q)$.  Write
\[
 \|A\|_1=\sum_{i,j}|A_{ij}|,\qquad
 \|A\|_{\max}=\max_{i,j}|A_{ij}|,
\]
for entrywise matrix norms.  Define the compact convex moment set
\begin{equation}
 \mathcal M(Q,K)=
 \left\{\int hh^{\mathsf T}\,d\nu(h):
        \operatorname{KL}(\nu\Vert P_Q)\leq K\right\}.
 \label{eq:mf-entropy-moment-set}
\end{equation}
Compactness follows from weak compactness of probability measures on
$[-1,1]^s$, lower semicontinuity of relative entropy, and continuity of
the bounded moment functions.  The reference covariance may be singular.

\subsection{Changing the reference covariance without increasing entropy}

\begin{lemma}[Entropy-ball covariance stability]
\label{lem:mf-entropy-stability}
Suppose $K>0$ and $\Delta=\|Q-Q'\|_{\max}\leq K$.
The Hausdorff distance between $\mathcal M(Q,K)$ and
$\mathcal M(Q',K)$ in maximum norm is at most
\begin{equation}
 2\sqrt{2K\Delta}+3\Delta\leq6\sqrt{K\Delta}.
 \label{eq:mf-entropy-stability}
\end{equation}
\end{lemma}
\begin{proof}
For $\|A\|_1\leq1$, put $U_A(h)=h^{\mathsf T}Ah$ and
$\psi_Q(\tau)=\log\mathbb E_{P_Q}e^{\tau U_A}$.
The support function of~\eqref{eq:mf-entropy-moment-set} is
\begin{equation}
 h_Q(A)=\inf_{\tau>0}\frac{K+\psi_Q(\tau)}\tau.
 \label{eq:mf-entropy-support}
\end{equation}
Here is an elementary justification valid for any reference law.
The entropy variational inequality gives the upper bound.
The tilted law $d\nu_\tau=e^{\tau U_A-\psi_Q(\tau)}dP_Q$ has
entropy $k(\tau)=\tau\psi_Q'(\tau)-\psi_Q(\tau)$.
This is continuous, nondecreasing, and tends to zero as
$\tau\downarrow0$.  If it reaches $K$, the tilted law proves
sharpness of the upper bound.  If it never exceeds $K$, these tilted
laws remain feasible, and their means tend to the essential supremum
of the bounded variable $U_A$ as $\tau\to\infty$.
The right side of~\eqref{eq:mf-entropy-support} has that same upper
limit.  These two cases establish equality, including degenerate
reference laws.

Restricting the infimum in~\eqref{eq:mf-entropy-support} to
$0<\tau\leq T$ increases it by at most $K/T$.
Indeed convexity and $\psi_Q(0)=0$ imply that
$\psi_Q(\tau)/\tau$ is nondecreasing.  For every $\tau>T$,
\[
 \frac{K+\psi_Q(T)}T
 \leq\frac KT+\frac{\psi_Q(\tau)}\tau
 \leq\frac KT+\frac{K+\psi_Q(\tau)}\tau.
\]
For $\tau\leq T$ the restricted infimum is already no larger than
that candidate, proving the assertion without assuming attainment.

Gaussian differentiation along the segment from $Q$ to $Q'$ gives
\begin{equation}
 |\psi_Q(\tau)-\psi_{Q'}(\tau)|
       \leq(3\tau+2\tau^2)\Delta.
 \label{eq:mf-logpartition-stability}
\end{equation}
To check the constants, set
$V(z)=\tau\tanh(z)^{\mathsf T}A\tanh(z)$.
The sums of the absolute first and second derivatives of $V$ are
bounded respectively by $2\tau$ and $6\tau$.
The covariance derivative of the log partition function is one half
of the tilted expectation of $\nabla^2V+\nabla V\nabla V^{\mathsf T}$,
paired with $Q'-Q$.  This proves~\eqref{eq:mf-logpartition-stability}.
Adding a small multiple of the identity and passing to the limit
justifies singular covariances; the functions being integrated are
bounded uniformly for fixed $\tau$.

It follows that, uniformly in $\|A\|_1\leq1$,
\[
 |h_Q(A)-h_{Q'}(A)|\leq K/T+(3+2T)\Delta.
\]
For $\Delta>0$, choose $T=\sqrt{K/(2\Delta)}$.
The case $\Delta=0$ is immediate.  This gives the first bound in
\eqref{eq:mf-entropy-stability}; its second uses $\Delta\leq K$.
Support-function comparison for compact convex sets, and duality of
entrywise $\ell^1$ and maximum norms, give the stated Hausdorff bound.
Equivalently, the support inequality implies that each moment set is
contained in the other plus the indicated maximum-norm ball, by the
finite-dimensional separating-hyperplane theorem.
\end{proof}

\begin{lemma}[A continuous choice of a nearby bounded-entropy tilt]
\label{lem:mf-regularized-tilt}
Let $P$ be a law on $[-1,1]^s$, let $A$ be a target second-moment
matrix, and suppose some $\nu$ satisfies
\[
 \operatorname{KL}(\nu\Vert P)\leq K,
 \qquad\left\|\int hh^{\mathsf T}d\nu-A\right\|_{\max}\leq\alpha,
 \qquad K,\alpha>0.
\]
Put $G=K/\alpha$ and $r=\alpha^2/(4K)$.
Maximize, over symmetric matrices of entrywise norm at most $G$,
\begin{equation}
 \langle\Gamma,A\rangle-
   \log\int e^{h^{\mathsf T}\Gamma h}\,dP(h)-r\|\Gamma\|_F^2.
 \label{eq:mf-regularized-dual}
\end{equation}
There is a unique maximizer.  Its quadratic exponential tilt $P_*$
satisfies
\[
 \operatorname{KL}(P_*\Vert P)\leq2K,
 \qquad\left\|\int hh^{\mathsf T}dP_*-A\right\|_{\max}\leq3\alpha.
\]
For $P=P_Q$, the maximizer is continuous in $(Q,A)$ and hence is a
Borel choice, including singular $Q$.
\end{lemma}
\begin{proof}
Write $J(\Gamma)=\langle\Gamma,A\rangle-\log\int e^{h^{\mathsf T}\Gamma h}dP$.
The entropy variational inequality with the given witness gives
$J(\Gamma)\leq K+G\alpha=2K$ throughout the feasible ball.
Strict concavity of~\eqref{eq:mf-regularized-dual} gives uniqueness.
Let $\Gamma_*$ be its maximizer and
$E=A-\int hh^{\mathsf T}dP_*$.
First-order optimality with the zero matrix gives
$\langle\Gamma_*,E\rangle\geq2r\|\Gamma_*\|_F^2\geq0$.
Thus
\[
 \operatorname{KL}(P_*\Vert P)
       =J(\Gamma_*)-\langle\Gamma_*,E\rangle\leq2K.
\]
First-order optimality with a feasible matrix maximizing its pairing
with $E$ also gives
\[
 G\|E\|_{\max}
 \leq\langle\Gamma_*,E\rangle+4rG^2
 \leq J(\Gamma_*)+4rG^2\leq3K.
\]
For symmetric $E$ the dual norm identity still holds on symmetric
matrices: an off-diagonal maximizing entry is represented by two
symmetric entries of half the weight.  Dividing by $G$ proves the
moment bound.
Finally the Gaussian integral in~\eqref{eq:mf-regularized-dual} is
jointly continuous in $Q,A,\Gamma$ on every fixed compact feasible
ball.  The unique maximizer is therefore continuous by compactness
and the elementary maximum theorem.
\end{proof}

\subsection{The improved finite cover}

\begin{lemma}[A Gaussian mean-field output cover]
\label{lem:mf-entropy-cover}
Fix $H\geq2$ and put $v_H=2^{H+1}$.  There are computable constants
$C_H\geq1$ such that, for $K\geq128$ and $0<a<1$, a deterministic
finite family $\mathcal C_{K,a}$ and Borel events $G_{N,K,a}$ satisfy
\begin{gather}
 \mu_N(G_{N,K,a}^c)\leq2(H+1)e^{-NK/16},\qquad
 N\|v\|_2^2\leq K\quad\hbox{on }G_{N,K,a},\notag\\
 \inf_{h\in\mathcal C_{K,a}}\sup_{x\in[-1,1]^d}|z_w(x)-h(x)|<a
       \quad\hbox{for every }w\in G_{N,K,a},\notag\\
 T_H(d,K,1/a)
    \leq C_H(d+1)K^{v_H-2}a^{-v_H}
                         \log\frac{C_H(d+1)K}{a},\notag\\
 \log|\mathcal C_{K,a}|
    \leq C_H(d+1)K^{v_H-1}a^{-v_H}
                         \log\frac{C_H(d+1)K}{a}.
 \label{eq:mf-entropy-cover-bounds}
\end{gather}
The assertions hold for every $N\geq T_H(d,K,1/a)$.
The family is independent of learner width and of the training sample.
Its centers may be chosen bounded in absolute value by $\sqrt K$.
\end{lemma}
\begin{proof}
Use an internal accuracy $\zeta=a/4$.
Require the norm event
\[
 \|W^1\|_F^2\leq NK,\qquad
 \|W^\ell\|_{\mathrm{op}}\leq\sqrt K\ (2\leq\ell\leq H),\qquad
 N\|v\|_2^2\leq K.
\]
The first-layer and readout chi-square bounds and Gaussian sphere-net
bound give failure at most $(H+1)e^{-NK/16}$ for $K\geq128$.
For example the internal estimate is
$2\exp[(2\log9-K/8)N]\leq e^{-NK/16}$.
On this event, $z_w$ has Lipschitz constant at most
$K^{(H+1)/2}$.
Choose a deterministic Euclidean input net of radius
$\zeta/(4K^{(H+1)/2})$ with cardinality
\[
 s\leq\left(1+\frac{8\sqrt d\,K^{(H+1)/2}}\zeta\right)^d,
\]
and a Borel nearest-net-point map.

Set the backward covariance tolerances
\begin{equation}
 \delta_H=\frac{\zeta^2}{8K},\qquad
 \delta_{\ell-1}=\frac{\delta_\ell^2}{2304K}
       \quad(2\leq\ell\leq H),\qquad
 G_\ell^{\rm orig}=\frac{8K}{\delta_\ell}.
 \label{eq:mf-entropy-budgets}
\end{equation}
For a reference feature law $P$ and empirical covariance $A$, define
\[
 \mathcal J_G(A;P)=\max_{\|\Gamma\|_1\leq G}
       \left[\langle\Gamma,A\rangle-
               \log\mathbb E_Pe^{h^{\mathsf T}\Gamma h}\right].
\]
The distribution-free covariance certificate gives
\begin{equation}
 \mathbb P\{\mathcal J_G>K\}
 \leq\exp\{q\log(2s^2+1)-NK/2\},\qquad
 q=\lceil G^2/K\rceil.
 \label{eq:mf-slack-certificate}
\end{equation}
For completeness, approximate every entrywise $\ell^1$-ball parameter
by the average of $q$ independent atoms from
$\{0,\pm Ge_1,\ldots,\pm Ge_{s^2}\}$.
The expected log-partition error is at most $G^2/(2q)\leq K/2$,
by Hoeffding's exponential-moment bound and Jensen.
The empirical dual is therefore at most $K/2$ plus the maximum over
at most $(2s^2+1)^q$ atomic averages.  Each averaged parameter has
exponential moment one for its empirical dual multiplied by $N$.
Markov and a union bound prove~\eqref{eq:mf-slack-certificate}.

Take $N$ at least
$4\lceil(G_1^{\rm orig})^2/K\rceil\log(2s^2+1)/K$.
The conditional Gaussian row law at every layer is exact, so
\eqref{eq:mf-slack-certificate} and a union bound show that all
$H$ layer certificates are at most $K$, except with probability
$He^{-NK/4}$.  Define $G_{N,K,a}$ by this certificate event and the
norm event.  Both are Borel.  The displayed exception bound follows.

Fix a draw in this event, and denote its actual activation covariances
on the net by $Q_1,\ldots,Q_H$.
The usual dual lift gives a law of entropy at most $K$ relative to
$P_{Q_{\ell-1}}$ whose moment is within
$K/G_\ell^{\rm orig}=\delta_\ell/8$ of $Q_\ell$.
At the first layer the reference is the fixed Gaussian input field.
To see the lift, a maximizing unregularized parameter satisfies
\[
 \mathcal J_G=\operatorname{KL}(P_*\Vert P)
                 +G\|A-\mathbb E_{P_*}hh^{\mathsf T}\|_{\max}.
\]
This is the first-order condition over the entrywise $\ell^1$ ball.

We next construct a finite sampler that approximates these covariances
while paying only a bounded entropy cost per row.  Its layer widths are
\begin{equation}
 J_\ell=\left\lceil32\delta_\ell^{-2}\log(8Hs^2)\right\rceil.
 \label{eq:mf-adaptive-widths}
\end{equation}
The auxiliary reference law has independent Gaussian incoming weights:
variance $1/d$ at the first layer and $1/J_{\ell-1}$ at layer
$\ell\geq2$, followed by readout variance $1/J_H^2$.
This one reference law depends only on $d,H,K,a$.

Suppose the preceding empirical covariance $\widehat Q_{\ell-1}$
is within $\delta_{\ell-1}$ of $Q_{\ell-1}$.
Lemma~\ref{lem:mf-entropy-stability}, followed by the preceding
certificate lift, supplies a law of entropy at most $K$ relative to
$P_{\widehat Q_{\ell-1}}$ whose moment is within
\[
 \delta_\ell/8+6\sqrt{K\delta_{\ell-1}}
       =\delta_\ell/4
\]
of $Q_\ell$.  At the first layer the reference is unchanged and the
same conclusion holds directly.
Apply Lemma~\ref{lem:mf-regularized-tilt} with
$\alpha=\delta_\ell/4$ and target $Q_\ell$.
Its tilt has entropy at most $2K$ per row and moment error at most
$3\delta_\ell/4$.
Generate the $J_\ell$ incoming Gaussian rows with this tilt,
independently given the past.  Their induced field has the required
reference covariance because the weights are divided by
$\sqrt{J_{\ell-1}}$.
Hoeffding and~\eqref{eq:mf-adaptive-widths} put their empirical moment
within $\delta_\ell/4$ of its expectation, except with conditional
probability at most $1/(4H)$.  Thus the resulting covariance is within
$\delta_\ell$ of $Q_\ell$.

After the first failure of a covariance condition, use the unmodified
Gaussian reference at all later layers.  This is stopping, not
conditioning, so no extra entropy cost is incurred.
The continuously selected tilts make this a Borel probability kernel.
The chain rule bounds its total hidden-weight relative entropy by
$2K\sum_\ell J_\ell$, and all covariance conditions succeed with
probability at least $3/4$.

Approximate-readout duality gives a vector $\beta$ with
\[
 \|\beta\|_1\leq K/\zeta,\qquad
 \|Q_H\beta-z_w|_{\rm net}\|_\infty\leq\zeta/4,
 \qquad \beta^{\mathsf T}Q_H\beta\leq K.
\]
For clarity, minimize $N\|u\|_2^2/2$ subject to approximating the
original output within $\zeta/4$ on the net.
The original readout is strictly feasible.  Convex optimality gives
$Nu=\Phi^{\mathsf T}\beta$, where $\Phi$ is the feature matrix.
Complementarity and Cauchy--Schwarz give
\[
 (\zeta/4)\|\beta\|_1=Nv^{\mathsf T}u-N\|u\|_2^2
 \leq\sqrt K\,r-r^2\leq K/4,\qquad r=\sqrt N\|u\|_2.
\]
The minimizer's norm is no larger than that of the feasible original
readout, proving the last inequality as well.

On successful hidden sampling, give readout coordinate $i$ mean
$v_i^*=(\beta\cdot h_i)/J_H$ and keep its reference variance
$1/J_H^2$; on failure use zero mean.
Since $\|\widehat Q_H-Q_H\|_{\max}\leq\delta_H$,
\[
 J_H\|v^*\|_2^2=\beta^{\mathsf T}\widehat Q_H\beta
       \leq K+\delta_H(K/\zeta)^2=9K/8.
\]
The conditional readout entropy is at most $9KJ_H/16$.
Thus the full comparison law $\nu_f$ for $f=z_w$ satisfies
\begin{equation}
 \operatorname{KL}(\nu_f\Vert\mu_{\rm aux})
 \leq B:=2K\sum_{\ell=1}^H J_\ell+9KJ_H/16.
 \label{eq:mf-auxiliary-entropy}
\end{equation}
Readout noise has variance at most $1/J_H$ at each net point.
Its maximum is at most $\zeta/8$ except with probability
$2s\exp(-J_H\zeta^2/128)\leq1/4$, by
\eqref{eq:mf-adaptive-widths} and~\eqref{eq:mf-entropy-budgets}.
Consequently, with comparison probability at least $1/2$, its output
evaluated at nearest-net points differs from $f$ everywhere by at most
\[
 \zeta/4\quad\hbox{(readout dual)}
 +\zeta/8\quad\hbox{(covariance)}
 +\zeta/8\quad\hbox{(readout noise)}
 +\zeta/4\quad\hbox{(input net)}
 =3\zeta/4<\zeta.
\]

Let $\rho_{\rm aux}$ be the distribution of these auxiliary
nearest-net-point functions under the single reference Gaussian law.
The distance from a fixed continuous $f$ to such a function is a
continuous function of its finitely many net values, so the balls
below are measurable.  Binary data processing for relative entropy gives
\[
 \rho_{\rm aux}\{h:\|h-f\|_\infty<\zeta\}
       \geq\exp[-2(B+\log2)].
\]
Indeed if a comparison law gives a set probability at least $1/2$,
its entropy is at least one half the negative logarithm of the
reference probability, minus $\log2$.
Every good output, at every original learner width, has this lower
bound under the same auxiliary measure.
Balls of radius $\zeta$ around functions separated by more than
$2\zeta$ are disjoint.  A maximal such packing therefore has at most
$\exp[2(B+\log2)]$ centers and covers the entire good output class
to accuracy $2\zeta=a/2<a$.
The packing centers may be selected from that class, so their absolute
values are at most $\sqrt K$.
This is a finite existence argument; no efficient construction of the
cover is claimed.

It remains to record the degrees.  Put $q_\ell=K/\delta_\ell$.
Then $q_H=8K^2/\zeta^2$ and
$q_{\ell-1}=2304q_\ell^2$.
At fixed depth,
\[
 \delta_1^{-2}\leq C_H K^{v_H-2}\zeta^{-v_H},
 \qquad v_H=2^{H+1}.
\]
The net logarithms are at most
$C_H(d+1)\log[C_H(d+1)K/\zeta]$.
The certificate width threshold is bounded by this logarithm times
$C_H\delta_1^{-2}$, giving the stated $T_H$ bound.
The widths $J_\ell$ decrease with $\ell$, and
\eqref{eq:mf-auxiliary-entropy} is at most $C_HKJ_1$.
This proves the cover-cardinality bound in
\eqref{eq:mf-entropy-cover-bounds}.  Replacing $\zeta$ by $a/4$
changes only constants depending on $H$.
\end{proof}

\section{The fixed mean-field prior: raw predictions at every width}
\label{sec:fixed-meanfield}

The next result keeps the prior prefactor equal to one.  It applies to
ordinary posterior predictions at finite width under mean-field scaling;
no additional prior-strength limit or sample-dependent prior prefactor is
used.  Its sample exponent is weaker than the strong-prior exponent, but
it holds in a regime where the typical internal Frobenius energy grows
with learner width.

There are $H\ge2$ hidden tanh layers, a scalar linear readout, and input
$\bar x=(1,x)$ with $x\in[-1,1]^d$, $d\ge1$.  The independent Gaussian
variances are
\[
 \operatorname{Var}(W^1_{jk})=1/d,\qquad
 \operatorname{Var}(W^\ell_{ji})=1/N\quad(2\le\ell\le H),
 \qquad \operatorname{Var}(v_j)=1/N^2.
\]
Write $\mu_N$ for this prior and use observation variance $n\lambda/N$.
The actual finite posterior is therefore
\begin{equation}\label{eq:mf-raw-posterior}
 \Pi_{N,S}(dw)=Z_{N,S}^{-1}
     e^{-N\widehat R_S(z_w)/(2\lambda)}\mu_N(dw).
\end{equation}
The IID data law is arbitrary on $[-1,1]^d\times[-1,1]$.

Put
\begin{equation}\label{eq:mf-improved-degrees}
 v_H=2^{H+1},\qquad D_H=v_H-1,\qquad A_H=v_H.
\end{equation}
The symbol $v_H$ is an exponent, not the network's readout vector.
For a width-$m$, weight-$B$ teacher $g$, let
\begin{gather}
 E=d(d+1)B^2+2Hm^2B^2,\qquad U=mB,\qquad
 \rho=\frac{\lambda}{4(U+1)},\notag\\
 b_0=2\sqrt{2(2\log9+\log(4(H+1)))},\qquad
 P_g=1+2mB(d+H+2)(2+3mB+b_0)^{H-1},\notag\\
 \eta_\rho=\min\{1/2,\rho/(2P_g)\},\qquad c_0=3/8,\notag\\
 A_g(\rho)=\frac E2+[d+2+(H-1)m]\log\frac1{c_0\eta_\rho}+\log2.
 \label{eq:mf-teacher-cost}
\end{gather}
The teacher-cylinder estimate gives
$\mu_N\{\|z_w-g\|_\infty<\rho\}\ge e^{-NA_g(\rho)}$
for every $N\ge m$.

\begin{theorem}[A finite mean-field oracle inequality for raw outputs]
\label{thm:fixed-meanfield-raw}
For every $H\ge2$ there are finite computable constants $C_H$ and
$\lambda_H>0$ with the following property.  Fix
$0<\lambda\le\min\{\epsilon,\lambda_H\}$, $\epsilon\le1$,
$0<\delta<1$, and $n\ge2$, and suppose
\begin{equation}\label{eq:mf-two-parameter-budget}
 n\ge C_H(d+1)^2
 \log^6\!\frac{2n(d+1)}{\lambda\epsilon\delta}
 \lambda^{-v_H}\epsilon^{-3}.
\end{equation}
For each fixed teacher $g$ with $B\ge1$, with probability at least
$1-\delta$, simultaneously at every learner width $N\ge m$,
\begin{align}
 \int R(z_w)\,\Pi_{N,S}(dw)
 \le\;&R(g)+2\lambda A_g(\rho)
       +C_H\left[\epsilon+\lambda\log\frac4\delta\right]
 \notag\\
 &+(U+1)\sqrt{\frac{2R(g)\log(4/\delta)}n}
       +\frac{2(U+1)^2\log(4/\delta)}{3n}.
 \label{eq:mf-raw-oracle}
\end{align}
The same bound holds for the risk of the ordinary posterior mean
$\int z_w\,d\Pi_{N,S}$, and for the usual clipped prediction
conventions.

For a fixed teacher, dimension, depth, and confidence level, setting
$\epsilon=\lambda$ gives the rate
\begin{equation}\label{eq:mf-raw-rate}
 \widetilde O\!\left(n^{-1/(2^{H+1}+3)}\right).
\end{equation}
Thus the exponents for two, three, and four hidden layers are
$1/11$, $1/19$, and $1/35$, respectively.
\end{theorem}

The constants in this theorem can grow rapidly with depth.  The theorem
improves finite mean-field coverage and the quantitative covariance
cover; it does not establish the conjectural exponent $1/H$ at this
fixed prior prefactor.

\subsection{A multiscale complexity with an exponential prior moment}

Set $K_0=128$ and let $K_j=K_0 2^j$, $0\le j\le J$, where $J$ is the least index with
$K_J\ge128/\lambda$.  Choose $\lambda_H$ small enough that
$K_J<256/\lambda$.  Fix a sufficiently small numerical $c_*>0$ and set
\[
 a_j=c_*(\epsilon+\lambda K_j),\qquad
 N_0=\left\lceil\max_j T_H(d,K_j,1/a_j)\right\rceil.
\]
For $N\ge N_0$, abbreviate the events in Lemma~\ref{lem:mf-entropy-cover}
by $G_j$ and define the sample-independent Borel quantity
\[
 q_N(w)=
 \begin{cases}
 K_0,&w\in\bigcap_jG_j,\\
 2\max\{K_j:w\notin G_j\},&\text{otherwise}.
 \end{cases}
\]
If $q_N=K_j$ for $j\le J$, then $w\in G_j$.  The only other
possibility is $q_N=2K_J$, in which case $w\notin G_J$.

\begin{lemma}[Entropy pays the multiscale complexity]
\label{lem:mf-multiscale-entropy}
With $B_H=K_0/128+\log(1+4(H+1))$,
\[
 \log\int e^{Nq_N/128}\,d\mu_N\le NB_H.
\]
Consequently the posterior in \eqref{eq:mf-raw-posterior} satisfies,
for every comparison distribution $Q$ of finite relative entropy,
\begin{equation}\label{eq:mf-complexity-comparison}
 \int\left[\widehat R_S(z_w)+\frac\lambda{64}q_N(w)\right]d\Pi
 \le\int\widehat R_S(z_w)\,dQ
       +\frac{2\lambda}{N}\KL(Q\Vert\mu_N)+2\lambda B_H.
\end{equation}
\end{lemma}
\begin{proof}
The cover failure bounds give
\[
 \int e^{Nq_N/128}\,d\mu_N
 \le e^{NK_0/128}+2(H+1)\sum_{j=0}^J e^{-3NK_j/64}
 \le e^{NK_0/128}+4(H+1)\le e^{NB_H}.
\]
The entropy variational inequality implies
$\KL(Q\Vert\mu_N)\ge N\int q_N\,dQ/128-NB_H$.
Apply this to the posterior and use its exact variational optimality
for $\int\widehat R_S\,dQ+(2\lambda/N)\KL(Q\Vert\mu_N)$.
\end{proof}

\subsection{Raw statistical deviations and posterior tails}

Clip each cover center, only for the proof, to
$[-\sqrt{K_j}-1,\sqrt{K_j}+1]$.  Its approximation to an output in
$G_j$ is preserved.  Adjoin its negative and zero if needed.  The
worst-case Rademacher complexity of this finite class is at most
\[
 (\sqrt{K_j}+1)
       \sqrt{2\log(2|\mathcal C_{K_j,a_j}|+1)/n}.
\]
Apply the smooth-loss inequality to all these finite classes with total
failure probability at most $\delta/4$.  If $q_N(w)=K_j$ and
$\widehat R_S(z_w)\le8$, an approximating center $h$ has empirical risk
less than $16$, provided $c_*$ was chosen small enough.  On the common
sample event,
\[
 R(h)-\widehat R_S(h)
 \le C\left[
 \sqrt{\frac{K_jM_j\log^3(2n)}n}
       +\frac{K_jM_j\log^3(2n)}n\right],
 \quad
 M_j=\log(2|\mathcal C_{K_j,a_j}|+1)
                 +\log\frac{4(J+1)}\delta.
\]
The norm inequalities
$|\sqrt{R(z_w)}-\sqrt{R(h)}|\le a_j$ and their empirical counterpart
then give
\begin{equation}\label{eq:mf-low-objective-deviation}
 R(z_w)\le\widehat R_S(z_w)+\frac\lambda{64}q_N(w)+C\epsilon,
\end{equation}
as soon as
\[
 n\ge C\log^3(2n)
          \max_j\frac{K_jM_j}{(\epsilon+\lambda K_j)^2}.
\]
Indeed, make the preceding center deviation a sufficiently small
multiple of $\epsilon+\lambda K_j$; its population risk then remains
bounded by a numerical constant, so transferring squared loss costs
only a constant times $a_j$.

The sample budget \eqref{eq:mf-two-parameter-budget} implies this
condition.  The relevant elementary maximization is
\[
 \sup_{x>0}\frac{x^{D_H+1}}
                     {(\epsilon+\lambda x)^{A_H+2}}
 \le C_H\lambda^{-v_H}\epsilon^{-2},
\]
since $D_H+1=A_H=v_H$ and $\epsilon\le1$.  The remaining factors are absorbed by the six
logarithms in that sample budget.

We next control the two regions excluded from
\eqref{eq:mf-low-objective-deviation}.  Put $r_N=N\|v\|^2$.
Under the prior, $r_N\sim\chi_N^2/N$, so
$\E r_N=1$, $\E r_N^2\le3$, and
\[
 R(z_w)\le2r_N+2,\qquad \int R(z_w)^2\,d\mu_N\le24.
\]
The centered readout gives $\int\widehat R_S(z_w)d\mu_N\le1+1/N$,
so Jensen yields
\begin{equation}\label{eq:mf-jensen-normalizer}
 Z_{N,S}\ge e^{-(N+1)/(2\lambda)}\ge e^{-N/\lambda}.
\end{equation}
The posterior risk contribution of $\widehat R_S(z_w)>8$ is at most
$2e^{-3N/\lambda}$.  For the top complexity state, Cauchy--Schwarz gives
\[
 \int R(z_w)\mathbf1_{G_J^{\mathsf c}}\,d\mu_N
 \le\sqrt{48(H+1)}e^{-NK_J/32}.
\]
Using $K_J\ge128/\lambda$ and
\eqref{eq:mf-jensen-normalizer}, its posterior contribution is at most
$\sqrt{48(H+1)}e^{-3N/\lambda}$.  Both tails are bounded by
$C_He^{-3/\lambda}$.

Integrate \eqref{eq:mf-low-objective-deviation} on its valid region,
extend the nonnegative empirical-risk-plus-complexity integral to the
whole posterior, and apply
\eqref{eq:mf-complexity-comparison}.  Take $Q$ to be the prior
conditioned on the teacher tube of radius $\rho$ in
\eqref{eq:mf-teacher-cost}.  Then
\[
 \int\widehat R_S(z_w)dQ\le\widehat R_S(g)+\lambda,
 \qquad \frac{2\lambda}{N}\KL(Q\Vert\mu_N)
             \le2\lambda A_g(\rho).
\]
This proves the theorem at every $N\ge\max\{m,N_0\}$ after the fixed
teacher's Bernstein bound.

\subsection{The raw small-width bridge}

For completeness, we give the bounded-loss argument that treats every
remaining width without changing the learner.  Set $M^2=128/\lambda$
and use the proof loss
$\ell_M=(\clip_{[-M,M]}z_w-Y)^2$, with range
$B=(M+1)^2\le512/\lambda$.
Put $\beta=N/(2\lambda)$ and $u=\beta B/n$.
For every fixed parameter, boundedness of this loss implies
\[
 \E_S\exp\left\{\frac nB(1-e^{-u})R(\clip_{[-M,M]}z_w)
                     -\beta\widehat R_S(\clip_{[-M,M]}z_w)\right\}
 \le1.
\]
Tonelli, Markov and entropy duality give, with sample failure at most
$\delta_N$,
\[
 \int R(\clip_{[-M,M]}z_w)d\Pi
 \le\frac u{1-e^{-u}}
 \left[\int\widehat R_S(\clip_{[-M,M]}z_w)d\Pi
       +\frac{\KL(\Pi\Vert\mu_N)}\beta
       +\frac{\log(1/\delta_N)}\beta\right].
\]
The posterior's raw variational objective is at most
$\int\widehat R_S(z_w)d\mu_N\le2$.  Clipping cannot increase empirical
loss.  Use $u/(1-e^{-u})\le1+u$ and the teacher comparison on the
unmultiplied objective.  When $u\le1$, this proves
\begin{align}
 \int R(\clip_{[-M,M]}z_w)d\Pi
 \le\;&\widehat R_S(g)+\lambda+2\lambda A_g(\rho)
 \notag\\
 &+\frac{4\lambda}{N}\log\frac1{\delta_N}
        +\frac{512N}{n\lambda^2}.
 \label{eq:mf-small-width-raw}
\end{align}
The clipping is solely a proof device; the likelihood still uses raw
outputs.  Since $|z_w(x)|\le\sqrt{r_N}$, a chi-square tail,
Cauchy--Schwarz, and \eqref{eq:mf-jensen-normalizer} give
\[
 \int R(z_w)\mathbf1_{\{r_N>M^2\}}d\Pi
                  \le\sqrt{24}e^{-7N/\lambda}.
\]
Adding this quantity restores the ordinary raw posterior risk in
\eqref{eq:mf-small-width-raw}.

Take $\delta_N=\delta/(4N^2)$ and intersect these events.  Their total
failure is less than $\delta/2$, and
$\log(1/\delta_N)/N\le\log(4/\delta)+1$.
By the threshold in \eqref{eq:mf-entropy-cover-bounds},
\[
 N_0\le C_H(d+1)\log\!\frac{2n(d+1)}{\lambda\epsilon\delta}
                 \lambda^{-(v_H-2)}\epsilon^{-2}.
\]
Here we used
$\sup_{x>0}x^{v_H-2}/(\epsilon+\lambda x)^{v_H}
\le C_H\lambda^{-(v_H-2)}\epsilon^{-2}$.
The sample budget \eqref{eq:mf-two-parameter-budget} therefore bounds
$512N_0/(n\lambda^2)$ by $C_H\epsilon$ and ensures $u\le1$ after
increasing its constant.  This covers every $m\le N<N_0$.

Intersect the large-width event, these small-width events, and the
fixed teacher's Bernstein event of failure $\delta/4$.  This proves
\eqref{eq:mf-raw-oracle} on one event for all learner widths.
The posterior has finite raw second moments because its density
relative to the Gaussian prior is bounded; Jensen proves the claim
for the ordinary posterior mean.

\subsection{Rates and growing comparator classes}

When $\epsilon=\lambda$, the sample budget has power $v_H+3$.
This proves \eqref{eq:mf-raw-rate}, with logarithmic factors included in
the calibration.

For width and weight caps $m,B\le n^\alpha$, choose
$\lambda=n^{-\theta}$ and $\epsilon=n^{-\eta}$ with
\[
 4\alpha<\theta<\frac1{v_H},\qquad
 0<\eta\le\theta,\qquad
 v_H\theta+3\eta<1.
\]
These choices exist for every
\begin{equation}\label{eq:mf-growing-class}
 \alpha<\frac1{4v_H}=2^{-(H+3)},
\end{equation}
by choosing $\eta>0$ sufficiently small.
The teacher energy is $O(n^{4\alpha})$, its cylinder logarithm is
smaller up to logarithms, and its Bernstein term tends to zero.
Choose a population-risk minimizer in the compact capped teacher class,
which contains zero and hence has minimum risk at most one.
The resulting coefficient-one oracle inequality has a vanishing
remainder uniformly over all learner widths above the cap.  For a fixed
weight cap and only width growing, replace $4$ by $2$ in
\eqref{eq:mf-growing-class}.

\section{An output cover uniform over stronger Gaussian scalings}
\label{sec:scaled-entropy-cover}

The entropy construction also applies when every Gaussian standard
deviation is multiplied by $\tau\in(0,1]$.  The first-layer, internal,
and readout variances are then respectively
$\tau^2/d$, $\tau^2/N$, and $\tau^2/N^2$.
For the scaling exponent $\gamma\geq1$, take
$\tau^2=N^{1-\gamma}$; its natural concentration speed is
$N/\tau^2=N^\gamma$.  We obtain a finite output cover common to all
$\tau$, although its cardinality is larger than the sharper
mean-field cover of Lemma~\ref{lem:mf-entropy-cover}.

\subsection{Normalized entropy and a uniform covariance bound}

For $0<\tau\leq1$, define
\[
 \mathcal M_\tau(Q,K)=
 \left\{\int hh^{\mathsf T}d\nu:
       \tau^2\operatorname{KL}(\nu\Vert P_{\tau^2Q})\leq K\right\}.
\]
The Hausdorff covariance estimate remains
\begin{equation}
 d_{\rm H,max}\bigl(\mathcal M_\tau(Q,K),
                   \mathcal M_\tau(Q',K)\bigr)
 \leq2\sqrt{2K\Delta}+3\Delta\leq6\sqrt{K\Delta},
 \qquad \Delta=\|Q-Q'\|_{\max}\leq K.
 \label{eq:scaled-entropy-stability}
\end{equation}
Indeed the support function in direction $A$, $\|A\|_1\leq1$, is
\[
 \inf_{u>0}\frac{K+\Psi_{\tau,Q}(u)}u,
 \qquad
 \Psi_{\tau,Q}(u)=\tau^2\log
   \mathbb E_{P_{\tau^2Q}}
            e^{(u/\tau^2)h^{\mathsf T}Ah}.
\]
The proof of~\eqref{eq:mf-entropy-support} applies after rescaling
the entropy constraint.  The covariance differentiation calculation
now gives
\[
 |\Psi_{\tau,Q}(u)-\Psi_{\tau,Q'}(u)|
       \leq\Delta(3u\tau^2+2u^2).
\]
Restricting $u$ to $(0,T]$ costs at most $K/T$, while division by $u$
bounds the covariance error by $(3+2T)\Delta$.  Optimizing $T$ proves
\eqref{eq:scaled-entropy-stability} exactly as before.  All assertions
include singular covariance matrices.

The continuous tilt-selection lemma has the same normalized version.
If a witness has normalized entropy at most $K$ and moment error at
most $\alpha$, maximize
\begin{equation}
 \langle\Gamma,A\rangle
 -\tau^2\log\mathbb E_{P_{\tau^2Q}}
                  e^{h^{\mathsf T}\Gamma h/\tau^2}
 -\frac{\alpha^2}{4K}\|\Gamma\|_F^2,
 \qquad \|\Gamma\|_1\leq K/\alpha.
 \label{eq:scaled-regularized-dual}
\end{equation}
Its unique maximizing tilt has normalized entropy at most $2K$ and
moment error at most $3\alpha$.  The proof of
Lemma~\ref{lem:mf-regularized-tilt} is unchanged, since the gradient
of the scaled log partition function is still the second moment.
The maximizing parameter is continuous in $(Q,A,\tau)$ for
$\tau>0$, giving the Borel kernels needed below.

\subsection{A common finite grid}

\begin{lemma}[Uniformly scaled Gaussian output cover]
\label{lem:scaled-mf-cover}
Fix $H\geq2$ and write
\[
 v_H=2^{H+1},\qquad D_H'=\tfrac32v_H-4,
 \qquad A_H'=\tfrac32v_H.
\]
There are computable constants $C_H\geq1$ such that for $K\geq128$
and $0<a<1$ a finite deterministic family $\mathcal D_{K,a}$,
independent of $\tau\in(0,1]$ and of learner width, satisfies
\begin{gather}
 T_H'(d,K,1/a)\leq
 C_H(d+1)K^{v_H-2}a^{-v_H}
                    \log\frac{C_H(d+1)K}{a},\notag\\
 \log|\mathcal D_{K,a}|\leq
 C_H(d+1)^2K^{D_H'}a^{-A_H'}
                    \log^3\frac{C_H(d+1)K}{a}.
 \label{eq:scaled-mf-cover-bounds}
\end{gather}
For every $N\geq T_H'(d,K,1/a)$ and every $\tau\in(0,1]$, there is
a Borel event $G_{N,\tau,K,a}$ such that
\begin{gather*}
 \mu_{N,\tau}(G_{N,\tau,K,a}^c)
       \leq2(H+1)\exp[-NK/(16\tau^2)],\qquad
 N\|v\|_2^2\leq K\quad\hbox{on }G_{N,\tau,K,a},\\
 \inf_{h\in\mathcal D_{K,a}}\sup_{x\in[-1,1]^d}|z_w(x)-h(x)|<a
       \quad(w\in G_{N,\tau,K,a}).
\end{gather*}
Every center may be chosen bounded by $\sqrt K$ in absolute value.
\end{lemma}
\begin{proof}
Set $\zeta=a/4$, and use the same input net and backward tolerances
as in~\eqref{eq:mf-entropy-budgets}:
\[
 s\leq\left(1+\frac{8\sqrt d\,K^{(H+1)/2}}\zeta\right)^d,
 \qquad \delta_H=\frac{\zeta^2}{8K},\quad
 \delta_{\ell-1}=\frac{\delta_\ell^2}{2304K},\quad
 G_\ell^{\rm orig}=8K/\delta_\ell.
\]
The norm event
\[
 \|W^1\|_F^2\leq NK,\qquad
 \|W^\ell\|_{\rm op}\leq\sqrt K\ (2\leq\ell\leq H),\qquad
 N\|v\|_2^2\leq K
\]
now fails with probability at most
$(H+1)\exp[-NK/(16\tau^2)]$.
For example the internal sphere-net estimate is
$2\exp[2N\log9-NK/(8\tau^2)]
 \leq\exp[-NK/(16\tau^2)]$.
The chi-square estimates for the first layer and readout have at least
the same strength.  On the norm event the output Lipschitz constant
is at most $K^{(H+1)/2}$, so the same net works for every $\tau$.

For an empirical activation covariance $A$ and a reference law $P$,
define the normalized dual certificate
\[
 \mathcal J_{G,\tau}(A;P)=
 \max_{\|\Gamma\|_1\leq G}
 \left[\langle\Gamma,A\rangle-
       \tau^2\log\mathbb E_P e^{h^{\mathsf T}\Gamma h/\tau^2}\right].
\]
For $N$ independent rows from $P$, the atomic sparsification argument
used in~\eqref{eq:mf-slack-certificate} gives
\begin{equation}
 \Pr\{\mathcal J_{G,\tau}>K\}
 \leq\exp\{q\log(2s^2+1)-NK/(2\tau^2)\},
 \qquad q=\lceil G^2/(K\tau^2)\rceil.
 \label{eq:scaled-certificate}
\end{equation}
Specifically, the expected scaled log-partition error in averaging
$q$ atoms is at most $G^2/(2q\tau^2)\leq K/2$.
For each fixed average its empirical dual, multiplied by
$N/\tau^2$, has exponential moment one.  These two facts prove
\eqref{eq:scaled-certificate} by the same finite union bound.
Since $G\geq K$ here, $\tau^2q\leq G^2/K+\tau^2\leq2G^2/K$.
Consequently the width requirement
\[
 N\geq8(G_1^{\rm orig})^2K^{-2}\log(2s^2+1)
\]
ensures failure at most $\exp[-NK/(4\tau^2)]$ at each layer,
uniformly in $\tau$.  Conditional row fields are exactly Gaussian
with covariance $\tau^2Q_{\ell-1}$, including the fixed first-layer
input field.  Intersect all the certificates and the norm event to
define $G_{N,\tau,K,a}$; the required probability bound follows.
The event is Borel because each certificate maximizes a continuous
function over a compact parameter set.

Fix a draw in this event.  Each normalized certificate at most $K$
supplies a law of normalized entropy at most $K$ whose moment is
within $\delta_\ell/8$ of the actual covariance $Q_\ell$.
The lift is the first-order identity
\[
 \mathcal J_{G,\tau}
 =\tau^2\operatorname{KL}(P_*\Vert P)
       +G\|A-\mathbb E_{P_*}hh^{\mathsf T}\|_{\max}.
\]
Use the deterministic unequal layer widths
\[
 J_\ell=\left\lceil32\delta_\ell^{-2}\log(8Hs^2)\right\rceil.
\]
Conditionally on previous sampled layers, take independent Gaussian
incoming rows of variance $\tau^2/d$ at the first layer and
$\tau^2/J_{\ell-1}$ thereafter, with the tilt selected by
\eqref{eq:scaled-regularized-dual} and $\alpha=\delta_\ell/4$.
Use the unmodified Gaussian rows after the first failure of a
covariance condition.  Just as in the proof of
Lemma~\ref{lem:mf-entropy-cover}, the normalized entropy stability
bound supplies a witness within $\delta_\ell/4$, the selected mean
is within $3\delta_\ell/4$, and Hoeffding contributes at most
$\delta_\ell/4$.  All sampled covariances are therefore within
$\delta_\ell$ of the actual ones with probability at least $3/4$.

The approximate-readout argument from that proof gives
$\|\beta\|_1\leq K/\zeta$ and
$\|Q_H\beta-z_w|_{\rm net}\|_\infty\leq\zeta/4$.
Set the sampled network's readout coefficient to
$v_i^*=(\beta\cdot h_i)/J_H$.  On successful covariance sampling,
its output on the net is within
$\zeta/4+(K/\zeta)\delta_H=3\zeta/8$ of the actual output.
Each readout coefficient has absolute value at most $K/(\zeta J_H)$.

We bound every incoming coefficient in one such realization,
uniformly in $\tau$.  There are at most
\[
 P_0=(d+1)J_1+\sum_{\ell=2}^H J_\ell J_{\ell-1}+J_H
\]
active coefficients, including the readout.  Every selected tilt
has parameter norm at most $4K/\delta_\ell\leq G_1^{\rm orig}$.
Its density relative to its incoming Gaussian row is therefore at
most $\exp(2G_1^{\rm orig}/\tau^2)$.  This also holds after stopping,
when that density is one.  Every actual Gaussian incoming coordinate
has variance at most $\tau^2$, since $d\geq1$ and
$J_{\ell-1}\geq1$.
Choose
\[
 B_0=\max\left\{1,\frac K{\zeta J_H},
       \sqrt{4G_1^{\rm orig}+2\log(8P_0)}\right\}.
\]
For every incoming coordinate, conditionally on all earlier layers,
its probability of absolute value above $B_0$ is at most
\[
 2\exp\left[-\frac{B_0^2-4G_1^{\rm orig}}{2\tau^2}\right]
       \leq2e^{-\log(8P_0)}=\frac1{4P_0}.
\]
A union bound gives simultaneous coefficient bounds with probability
at least $3/4$.  Intersecting with covariance success therefore gives
a deterministic realization with all these properties.  Importantly,
the bounds $J_\ell,B_0$ and the input net do not depend on $\tau$.

It remains to count bounded networks, rather than to use the
mean-field mass-packing bound, whose entropy cost would now contain
$1/\tau^2$.  Put $M=\max_\ell J_\ell$ and
\[
 \Lambda=(d+1+(H+1)M)[1+M(B_0+1)]^H.
\]
Round each of the at most $P_0$ actual coefficients to a grid of mesh
$\zeta/(8\Lambda)$, keeping the fixed unequal layer widths.
The layerwise coordinate-error induction, using $|\tanh|\leq1$ and
the one-Lipschitz property, bounds the resulting output perturbation
by $\zeta/8$ on the cube.  The grid contains at most
\[
 \exp\left[P_0\log\left(3+\frac{16B_0\Lambda}{\zeta}\right)\right]
\]
networks.  Evaluate their outputs at the fixed nearest input-net
point, and clip the values to $[-\sqrt K,\sqrt K]$.
Clipping cannot increase error to a good output, which already has
this range.  The total error is at most
$3\zeta/8+\zeta/8+\zeta/4=3\zeta/4<a$.
These functions form a single deterministic cover for every $\tau$.

For the stated degrees, write $q_\ell=K/\delta_\ell$.
The recursion $q_H=8K^2/\zeta^2$,
$q_{\ell-1}=2304q_\ell^2$ gives
\[
 J_\ell\leq C_H(d+1)
   K^{2^{H-\ell+2}-2}\zeta^{-2^{H-\ell+2}}
                \log\frac{C_H(d+1)K}{\zeta}.
\]
The uniform width threshold is a constant times
$\delta_1^{-2}\log(2s^2+1)$, with degrees $v_H-2$ and $v_H$.
Among the adjacent products in $P_0$, $J_1J_2$ is largest and has
degrees $3v_H/2-4$ and $3v_H/2$; these also dominate the first-layer
and readout counts for $H\geq2$ and $K\geq128$.
The rounding logarithm is bounded by
$C_H\log[C_H(d+1)K/\zeta]$.  This proves
\eqref{eq:scaled-mf-cover-bounds}, after absorbing $\zeta=a/4$
into constants depending only on $H$.
\end{proof}

\section{A sharper cover when the Gaussian scale is bounded below}
\label{sec:truncated-scale-mass-cover}

The finite-grid cover from Lemma~\ref{lem:scaled-mf-cover} is uniform
even as the Gaussian standard deviation tends to zero.  When its scale
is bounded below, the more economical entropy packing argument can
also be made uniform.  We use an auxiliary probability measure only
to count separated output functions.  This does not change the
learner's Gaussian prior, likelihood, or fixed architecture.

\begin{lemma}[A common entropy cover above a scale threshold]
\label{lem:truncated-scale-cover}
Fix $H\geq2$, put $v_H=2^{H+1}$, and let
$K\geq128$, $0<a<1$, and $0<\kappa\leq1$.
There are computable constants $C_H$ and a finite deterministic family
$\mathcal E_{K,a,\kappa}$, common to every scale
$\tau\in[\kappa,1]$ and every learner width, such that
\begin{gather}
 T_H''(d,K,1/a)\leq
 C_H(d+1)K^{v_H-2}a^{-v_H}
                      \log\frac{C_H(d+1)K}{a},\notag\\
 \log|\mathcal E_{K,a,\kappa}|\leq
 C_H\kappa^{-2}(d+1)K^{v_H-1}a^{-v_H}
                      \log\frac{C_H(d+1)K}{a\kappa}.
 \label{eq:truncated-scale-cover-bounds}
\end{gather}
For every $N\geq T_H''(d,K,1/a)$ and every
$\tau\in[\kappa,1]$, the Borel good event from
Lemma~\ref{lem:scaled-mf-cover} obeys
\begin{gather*}
 \mu_{N,\tau}(G_{N,\tau,K,a}^c)
       \leq2(H+1)e^{-NK/(16\tau^2)},\qquad
 N\|v\|_2^2\leq K\quad\hbox{on }G_{N,\tau,K,a},\\
 \inf_{h\in\mathcal E_{K,a,\kappa}}
          \sup_{x\in[-1,1]^d}|z_w(x)-h(x)|<a
             \quad(w\in G_{N,\tau,K,a}).
\end{gather*}
The centers can be chosen bounded in absolute value by $\sqrt K$.
\end{lemma}
\begin{proof}
Retain the net, $\zeta=a/4$, the tolerances $\delta_\ell$, and widths
$J_\ell$ from the preceding two cover constructions.  Consider the
single fixed auxiliary architecture with these unequal layer widths.
For a scale $\tau$, its reference measure $\mu_\tau^{\rm aux}$ has
independent centered Gaussian coordinates of variances
$\tau^2/d$ in the first layer, $\tau^2/J_{\ell-1}$ in internal layer
$\ell$, and $\tau^2/J_H^2$ in the readout.
Its number of coordinates, denoted $P_{\rm aux}$, satisfies
\[
 P_{\rm aux}\leq(d+1)J_1+
                   \sum_{\ell=2}^H J_\ell J_{\ell-1}+J_H.
\]
This architecture is used only for the counting argument.

Fix an actual good output $f$ at any learner width and any scale
$\tau\in[\kappa,1]$.  The normalized adaptive tilts from
\eqref{eq:scaled-regularized-dual}, stopped after the first covariance
failure, have conditional row relative entropy at most $2K/\tau^2$.
On successful hidden sampling use the readout mean
$v_i^*=(\beta\cdot h_i)/J_H$, and on failure use zero mean, retaining
reference readout variance $\tau^2/J_H^2$ in both cases.
As in Lemma~\ref{lem:mf-entropy-cover},
$J_H\|v^*\|_2^2\leq9K/8$, so the readout entropy is at most
$9KJ_H/(16\tau^2)$.  The full comparison law $\nu_{f,\tau}$ therefore
satisfies
\begin{equation}
 \operatorname{KL}(\nu_{f,\tau}\Vert\mu_\tau^{\rm aux})
       \leq B/\tau^2,
 \qquad B=2K\sum_{\ell=1}^H J_\ell+9KJ_H/16.
 \label{eq:scaled-auxiliary-entropy}
\end{equation}
The covariance sampling succeeds with probability at least $3/4$.
Readout noise at each net point has variance at most
$\tau^2/J_H\leq1/J_H$, so its maximum is at most $\zeta/8$ except
with probability at most $1/4$.  Thus $\nu_{f,\tau}$ gives probability
at least $1/2$ to auxiliary nearest-net-point outputs within $\zeta$
of $f$ uniformly on the input cube.  These are exactly the same
approximation bounds used for the mean-field entropy cover.

To obtain one counting measure for all $\tau$, set
$r=1+1/P_{\rm aux}$ and let
\[
 \tau_j=\kappa r^j\quad(0\leq j\leq J_*),
 \qquad J_* = \min\{j\geq0:\kappa r^j\geq1\}.
\]
Let $M_*=J_*+1$ and use the equally weighted auxiliary measure
\[
 \overline\mu^{\rm aux}=\frac1{M_*}
                   \sum_{j=0}^{J_*}\mu_{\tau_j}^{\rm aux}.
\]
The finite-dimensional Gaussian covariance matrices of these measures
are all scalar multiples of the same positive definite diagonal matrix.
For every $\tau\in[\kappa,1]$, a grid point satisfies
$\tau\leq\tau_j\leq r\tau$.  At every weight vector,
\[
 \frac{d\mu_{\tau_j}^{\rm aux}}{d\mu_\tau^{\rm aux}}
 =\left(\frac\tau{\tau_j}\right)^{P_{\rm aux}}
   \exp\left[\frac{Q_{
       \rm aux}}2(\tau^{-2}-\tau_j^{-2})\right]
 \geq r^{-P_{\rm aux}}\geq e^{-1},
\]
where $Q_{\rm aux}\geq0$ is the reference diagonal quadratic form.
Consequently
\begin{equation}
 \operatorname{KL}(\nu_{f,\tau}\Vert\overline\mu^{\rm aux})
 \leq B/\tau^2+1+\log M_*
 \leq B/\kappa^2+1+\log M_*.
 \label{eq:common-scale-entropy}
\end{equation}
No scale is mixed in the actual learner.  The auxiliary measure in
\eqref{eq:common-scale-entropy} is solely a device proving that each
of many disjoint function balls has nonnegligible mass under one
probability measure.

Push $\overline\mu^{\rm aux}$ forward to auxiliary outputs evaluated
at the fixed nearest-net-point map.  Binary relative-entropy data
processing and the comparison probability at least $1/2$ show that
every $\zeta$-ball around a good actual output has reference mass at
least
\[
 \exp\{-2[B/\kappa^2+1+\log M_*+\log2]\}.
\]
A maximal packing with pairwise distances greater than $2\zeta$
therefore has at most the reciprocal of this mass in cardinality and
is a $2\zeta$-cover of all good outputs, simultaneously over learner
width and $\tau\in[\kappa,1]$.  Choose its centers from the good
outputs themselves, so they have magnitude at most $\sqrt K$.

Finally, $\log(1+1/P_{\rm aux})\geq1/(2P_{\rm aux})$ implies
\[
 M_*\leq2+2P_{\rm aux}\log(1/\kappa).
\]
The degrees of $P_{\rm aux}$ are those in
Lemma~\ref{lem:scaled-mf-cover}, while
$B\leq C_HKJ_1$ has degrees $v_H-1$ and $v_H$.
Thus $\log M_*$ is absorbed by the logarithm in
\eqref{eq:truncated-scale-cover-bounds}.  The width threshold remains
the same uniform normalized-certificate threshold, independent of
$\kappa$.  This proves the asserted bounds.
\end{proof}

\subsection{One small-width event for every Gaussian scale}
\label{sec:simultaneous-scale-pacbayes}

The small-width event can be made simultaneous over an interval of
prior scales.  A finite mixture is used only as the reference measure
in the proof of a PAC--Bayes inequality.  Each learner still has its
specified single dense Gaussian prior.

For width $N$ let $\mu_{N,\tau}$ have first-layer, internal, and readout
variances $\tau^2/d$, $\tau^2/N$, and $\tau^2/N^2$, respectively.
Let $\Pi_{N,\tau,S}$ be its posterior with raw Gaussian likelihood
variance $n\lambda\tau^2/N$.  Its inverse empirical-risk temperature is
$\beta_\tau=N/(2\lambda\tau^2)$.

\begin{lemma}[Simultaneous-scale raw PAC--Bayes bound]
\label{lem:simultaneous-scale-small-width}
Fix $N\geq1$, $0<\kappa\leq1$, $0<\lambda\leq1$, and
$0<\delta'<1$.  Put
\[
 P_N=N(d+2)+(H-1)N^2,\qquad
 M_N=1+\left\lceil
       \frac{\log(1/\kappa)}{\log(1+1/P_N)}\right\rceil.
\]
Assume $256N/(n\lambda^2\kappa^2)\leq1$.
With sample probability at least $1-\delta'$, simultaneously for every
$\tau\in[\kappa,1]$ and every fixed eligible teacher $g$,
\begin{align}
 \int R(z_w)\,\Pi_{N,\tau,S}(dw)
 &\leq\widehat R_S(g)+\lambda+2\lambda\mathcal A_g(\lambda)
       +\frac{512N}{n\lambda^2\kappa^2}\notag\\
 &\quad+\frac{4\lambda\kappa^2}{N}
       [1+\log M_N+\log(1/\delta')]
       +\sqrt{24}e^{-7N/\lambda}.
 \label{eq:simultaneous-scale-small-bound}
\end{align}
Here $\mathcal A_g(\lambda)$ is the teacher-cylinder quantity
in~\eqref{eq:raw-small-Ag}.  The same bound holds for each ordinary
posterior mean.  The event does not depend on $\tau$ or on the teacher.
\end{lemma}

\begin{proof}
Let $r_N=1+1/P_N$ and form the grid
$\tau_j=\min\{1,\kappa r_N^j\}$ for $j=0,\ldots,M_N-1$.
The last point is one.  For every $\tau\in[\kappa,1]$ there is a grid
point with $\tau\leq\tau_j\leq r_N\tau$.
As all Gaussian coordinate variances scale together, their densities
satisfy the pointwise comparison
\[
 \frac{d\mu_{N,\tau_j}}{d\mu_{N,\tau}}(w)
 =\left(\frac\tau{\tau_j}\right)^{P_N}
   \exp\left[\frac{\mathcal Q_N(w)}2
                  (\tau^{-2}-\tau_j^{-2})\right]
 \geq r_N^{-P_N}\geq e^{-1}.
\]
Here $\mathcal Q_N$ is the nonnegative quadratic form for the
unit-scale prior; its exact expression is unnecessary in this comparison.
For the single auxiliary reference
$\overline\mu_N=M_N^{-1}\sum_j\mu_{N,\tau_j}$, it follows that
\begin{equation}
 \operatorname{KL}(Q\Vert\overline\mu_N)
 \leq\operatorname{KL}(Q\Vert\mu_{N,\tau})+1+\log M_N
 \label{eq:scale-grid-domination}
\end{equation}
whenever the right side is finite.

Set $M^2=128/\lambda$ and use the proof loss of
$z_w^{[M]}=\operatorname{clip}_{[-M,M]}z_w$.
Its range is at most $b=(M+1)^2\leq512/\lambda$.
Apply the bounded-loss Catoni argument from
Lemma~\ref{lem:raw-small-width} once, with reference
$\overline\mu_N$ and the fixed parameter
\[
 \beta_* =\frac{N}{2\lambda\kappa^2},\qquad
 u_* =\frac{\beta_*b}{n}\leq1.
\]
On one event, every probability law $Q$ of finite relative entropy
obeys that PAC--Bayes inequality.  In particular it applies to every
$\Pi_{N,\tau,S}$ without a union over $\tau$.

For such a posterior, define
\[
 A_\tau=\int\widehat R_S(z_w^{[M]})\,d\Pi_{N,\tau,S}
      +\frac{\operatorname{KL}(\Pi_{N,\tau,S}\Vert\mu_{N,\tau})}
                   {\beta_*}.
\]
Since $\beta_*\geq\beta_\tau$, relative entropy is nonnegative, and
clipping cannot increase empirical loss,
\[
 0\leq A_\tau
 \leq\int\widehat R_S(z_w)\,d\Pi_{N,\tau,S}
       +\frac{\operatorname{KL}(\Pi_{N,\tau,S}\Vert\mu_{N,\tau})}
                    {\beta_\tau}
 \leq\int\widehat R_S(z_w)\,d\mu_{N,\tau}\leq2.
\]
The penultimate inequality compares the true posterior objective with
its own prior.  The same objective is at most
$\widehat R_S(g)+\lambda+2\lambda\mathcal A_g(\lambda)$ by the
teacher comparison, because $\tau\leq1$.
Use~\eqref{eq:scale-grid-domination},
$u_* /(1-e^{-u_*})\leq1+u_*$, and $A_\tau\leq2$ in the PAC--Bayes
bound.  This gives
\begin{align*}
 \int R(z_w^{[M]})\,d\Pi_{N,\tau,S}
 &\leq A_\tau+2u_*
        +\frac{2}{\beta_*}[1+\log M_N+\log(1/\delta')]\\
 &\leq\widehat R_S(g)+\lambda+2\lambda\mathcal A_g(\lambda)
       +\frac{512N}{n\lambda^2\kappa^2}
       +\frac{4\lambda\kappa^2}{N}
                       [1+\log M_N+\log(1/\delta')].
\end{align*}
The weighted readout tail proved in Lemma~\ref{lem:raw-small-width}
is deterministic in the sample and holds for each $\tau$ with
$\sigma^2=\tau^2/N$.  It is at most
$\sqrt{24}\exp[-7N/(\lambda\tau^2)]\leq\sqrt{24}e^{-7N/\lambda}$,
so adding it proves~\eqref{eq:simultaneous-scale-small-bound}.
All pointwise scale comparisons are valid on the same PAC--Bayes
event.  Jensen proves the posterior-mean assertion.
\end{proof}

Taking $\delta'_N=\delta/(4N^2)$ and intersecting the countably many
width events gives total failure below $\delta/2$.
The extra scale-grid term remains uniformly bounded.  Indeed, with
$D_0=d+H+1$ and $L=\log(1/\kappa)$,
\[
 M_N\leq2+(P_N+1)L
       \leq N^2[2+(D_0+1)L].
\]
The inequalities $4\log N/N\leq2$ and
$\kappa^2\log(1+L)\leq1$ give
\begin{equation}
 \frac{4\lambda\kappa^2}{N}
       [1+\log M_N+\log(4N^2/\delta)]
 \leq4\lambda\left[4+\log\frac{4(D_0+3)}\delta\right].
 \label{eq:simultaneous-scale-confidence}
\end{equation}
Thus the simultaneous-scale event has the same sample power as the
pointwise-scale small-width argument.  A fixed-teacher Bernstein event
can be shared with the large-width and spectral regimes.

\section{One fixed-prefactor oracle bound for every exponent at least one}
\label{sec:fixed-all-gamma-raw}

The two approximation arguments cover complementary Gaussian scales.
The covariance cover works when the internal Gaussian scale is bounded
below; spectral truncation works when that scale is small.  Combining
them gives a rate independent of the fixed exponent $\gamma\geq1$,
without modifying the Gaussian prior's prefactor.

Use $H\geq2$ hidden tanh layers, inputs in $[-1,1]^d$, and bounded
IID labels $|Y|\leq1$.  For each fixed $\gamma\geq1$, the actual
independent centered Gaussian variances are
\[
 \operatorname{Var}(W^1_{jk})=N^{1-\gamma}/d,\qquad
 \operatorname{Var}(W^\ell_{ji})=N^{-\gamma},\qquad
 \operatorname{Var}(v_j)=N^{-\gamma-1}.
\]
The Gaussian likelihood variance is $n\lambda N^{-\gamma}$.
Thus the posterior is exactly
\[
 \Pi_{N,\gamma,S}(dw)=Z^{-1}
        e^{-N^\gamma\widehat R_S(z_w)/(2\lambda)}\mu_{N,\gamma}(dw).
\]
Write $A_g(\rho),E,U$ for the teacher quantities in
\eqref{eq:mf-teacher-cost}; in particular
$E=d(d+1)B^2+2Hm^2B^2$, $U=mB$, and
$\rho=\lambda/[4(U+1)]$.

\begin{theorem}[A raw oracle bound at every mean-field-or-stronger exponent]
\label{thm:fixed-all-gamma-raw}
Let
\[
 v_H=2^{H+1},\qquad b_H=\max\{1,(H-1)/2\},\qquad
 a_H=v_H+2b_H.
\]
There are finite computable constants $C_H$ and $\lambda_H>0$, depending
only on $H$, with the following property.  Suppose $n\geq2$,
$0<\lambda\leq\min\{\epsilon,\lambda_H\}$, $\epsilon\leq1$,
$0<\delta<1$, and
\begin{equation}
 n\geq C_H(d+1)^2
    \log^6\!\frac{2n(d+1)}{\lambda\epsilon\delta}\,
                 \lambda^{-a_H}\epsilon^{-3}.
 \label{eq:fixed-all-gamma-budget}
\end{equation}
For each fixed width-$m$, weight-$B$ teacher $g$ with $B\geq1$, with
probability at least $1-\delta$, simultaneously for every
$\gamma\geq1$ and every learner width $N\geq m$,
\begin{align}
 \int R(z_w)\,\Pi_{N,\gamma,S}(dw)
 \leq\;&R(g)+2\lambda A_g(\rho)
            +C_H\left[\epsilon+\lambda\log\frac{8(d+H+4)}\delta\right]
 \notag\\
 &+(U+1)\sqrt{\frac{2R(g)\log(8/\delta)}n}
        +\frac{2(U+1)^2\log(8/\delta)}{3n}.
 \label{eq:fixed-all-gamma-oracle}
\end{align}
The same bound holds for the ordinary raw-output posterior mean and
for the usual clipped prediction conventions.
The constants, calibration, and probability event are common to all
the stated widths and prior exponents.

For a fixed teacher, dimension, depth, and confidence level, setting
$\epsilon=\lambda$ gives the excess rate
\begin{equation}
 \widetilde O\!\left(n^{-1/(2^{H+1}+2\max\{1,(H-1)/2\}+3)}\right).
 \label{eq:fixed-all-gamma-rate}
\end{equation}
For two, three, and four hidden layers these exponents are
$1/13$, $1/21$, and $1/38$.  For every
\begin{equation}
 \alpha<\frac1{4a_H},
 \label{eq:fixed-all-gamma-growing}
\end{equation}
the posterior also competes, with vanishing additive excess, with the
best population risk in the class of teachers whose width and weight
caps are both at most $n^\alpha$.  At a fixed weight cap the sufficient
width exponent is twice as large.
\end{theorem}

\begin{proof}
Put $\tau_N=N^{(1-\gamma)/2}$ and $\omega_N=N^\gamma=N/\tau_N^2$.
Choose a sufficiently small constant $c_H>0$, depending only on $H$,
and introduce the proof threshold
\[
 \kappa=c_H\lambda^{b_H}.
\]
This threshold only selects the argument used for a given width; it
does not change any weight variance.

First consider widths with $\tau_N<\kappa$.
The sample condition \eqref{eq:fixed-all-gamma-budget} implies the
spectral statistical calibrations \eqref{eq:raw-all-lambda}, with
failure parameter $\delta/4$.  Indeed $a_H\geq H$, while its logarithm
and $(d+1)^2$ factors dominate those in the displayed calibrations.
It also makes
\[
 r_n+\frac{512C_0^2Tu_n}{\lambda}\leq C_H\epsilon,
 \qquad M_n\leq\overline M_H:=7T
\]
for a finite constant $\overline M_H$ depending only on $H$;
here $T=64(1+c_H^{\rm eff})$ and $c_H^{\rm eff}$ is the constant
in \eqref{eq:effective-mgf-constant}.
The dimension-free bound follows because the effective energy omits
the first layer, whose weights are unrestricted in the capacity bound.
Taking the threshold constant $c_H$ no larger than
\[
 \left[768\sqrt{\overline M_H}(H-1)2^{(H-1)/2}
       T^{\max\{(H-3)/2,0\}}\right]^{-1}
\]
ensures \eqref{eq:raw-all-prefactor}.  Theorem~\ref{thm:raw-all-scales}
then proves \eqref{eq:fixed-all-gamma-oracle} at all these widths on
one common event.  Its teacher contribution is at most
$2\lambda A_g(\rho)+C_H\lambda$, since $\tau_N\leq1$;
its remaining spectral transfer term is $O_H(\kappa)\leq O_H(\lambda)$.

It remains to treat $\tau_N\geq\kappa$.  Set $K_0=128$,
$K_j=K_0 2^j$ up to the least $J$ for which $K_J\geq128/\lambda$,
and let $a_j=c_*(\epsilon+\lambda K_j)$ for the same sufficiently
small numerical constant as in the mean-field proof.
Apply Lemma~\ref{lem:truncated-scale-cover} to these levels.
Its width threshold has no $\kappa$ dependence, so
\begin{equation}
 N_0:=\left\lceil\max_j T_H''(d,K_j,1/a_j)\right\rceil
 \leq C_H(d+1)\log\!\frac{2n(d+1)}{\lambda\epsilon\delta}\,
                   \lambda^{-(v_H-2)}\epsilon^{-2}.
 \label{eq:fixed-all-gamma-width-threshold}
\end{equation}
At $N\geq N_0$, define $q_{N,\gamma}$ to be twice the largest failed
cover level, or $K_0$ when no level fails.  The normalized Gaussian
failure probabilities give exactly
\[
 \log\int e^{\omega_Nq_{N,\gamma}/128}\,d\mu_{N,\gamma}
       \leq\omega_N B_H.
\]
Thus the posterior variational comparison pays
$(\lambda/64)q_{N,\gamma}$ as in
\eqref{eq:mf-complexity-comparison}.

The cover family is common to every $\tau\in[\kappa,1]$ and every
width.  Its cardinality bound differs from the mean-field bound only
by a factor $\kappa^{-2}$ and the logarithm involving $1/\kappa$.
The finite-class smooth-loss calculation from the mean-field proof
therefore applies on one common sample event whenever
\begin{equation}
 n\geq C_H(d+1)^2
    \log^6\!\frac{2n(d+1)}{\lambda\epsilon\delta}\,
                  \kappa^{-2}\lambda^{-v_H}\epsilon^{-2}.
 \label{eq:fixed-all-gamma-large-width-budget}
\end{equation}
In detail, the relevant maximum is bounded by
\[
 \kappa^{-2}\sup_{x>0}
       \frac{x^{v_H}}{(\epsilon+\lambda x)^{v_H+2}}
       \leq C_H\kappa^{-2}\lambda^{-v_H}\epsilon^{-2}.
\]
The log factor involving $\kappa=c_H\lambda^{b_H}$ is absorbed into
$C_H$ times the displayed logarithm.
On empirical risk at most eight and below the top failed complexity
level, the raw risk deviation is at most
$\lambda q_{N,\gamma}/64+C_H\epsilon$.

The posterior tail estimates from the mean-field proof also remain
valid.  Under the scaled prior,
$r_N=N\|v\|^2=N^{-\gamma}\chi_N^2$ has first and second moments
at most one and three.  Hence $\int R(z_w)^2d\mu\leq24$, and the
centered readout implies $Z\geq e^{-\omega_N/\lambda}$.
Both omitted weighted posterior tails are at most
$C_He^{-3\omega_N/\lambda}\leq C_He^{-3/\lambda}$.
Integration of the raw deviation, followed by the teacher-cylinder
comparison, proves the desired bound for all widths with
$\tau_N\geq\kappa$ and $N\geq N_0$.

For the remaining widths use
Lemma~\ref{lem:simultaneous-scale-small-width}, with failure level
$\delta/(8N^2)$.  At each width its one event covers every
$\tau\in[\kappa,1]$, hence every eligible $\tau_N$ without an
uncountable union.  Its sample remainder is at most
\[
 \frac{512N}{n\lambda^2\kappa^2}
     \leq\frac{512N_0\kappa^{-2}}{n\lambda^2}.
\]
Using \eqref{eq:fixed-all-gamma-width-threshold}, this is at most
$C_H\epsilon$ whenever
\[
 n\geq C_H(d+1)
       \log\!\frac{2n(d+1)}{\lambda\epsilon\delta}\,
                     \kappa^{-2}\lambda^{-v_H}\epsilon^{-3}.
\]
Increasing its constant also ensures the lemma's calibration.
This requirement and \eqref{eq:fixed-all-gamma-large-width-budget}
follow from \eqref{eq:fixed-all-gamma-budget}, because
$\kappa^{-2}=c_H^{-2}\lambda^{-2b_H}$.
The additional scale-grid entropy and confidence costs are uniform
in width and scale by \eqref{eq:simultaneous-scale-confidence}:
\[
 \frac{4\lambda\kappa^2}{N}
    [1+\log M_N+\log(8N^2/\delta)]
 \leq4\lambda\left[4+\log\frac{8(d+H+4)}\delta\right].
\]
The auxiliary scale mixture used in that lemma only constructs its
common sample event.  The posteriors in this theorem retain their
individual Gaussian priors.
The small-width exceptional probabilities sum to less than $\delta/4$.
Together with the two common statistical events and one fixed-teacher
Bernstein event, the total failure probability is at most $\delta$.
Every event is now common to all the asserted exponents and widths.
Jensen and posterior square integrability prove the assertion for the
ordinary mean.

Finally $\epsilon=\lambda$ gives sample power $a_H+3$.
For growing teacher caps take $\lambda=n^{-\theta}$ and
$\epsilon=n^{-\eta}$, where
\[
 4\alpha<\theta<1/a_H,\qquad
 0<\eta\leq\theta,\qquad a_H\theta+3\eta<1.
\]
These choices exist under \eqref{eq:fixed-all-gamma-growing}.
They make $\lambda E$ and every other oracle remainder vanish for the
population-risk minimizer of the compact capped class.  The minimizer
is fixed before sampling and has risk at most one, because the zero
network belongs to the class.  With weights fixed, replace $4\alpha$
by $2\alpha$.
\end{proof}

\subsection{A two-hidden-layer refinement for growing teacher classes}
\label{sec:two-layer-scale-refinement}

For two hidden layers the scale threshold can depend on both accuracy
parameters, improving the growing-comparator exponent while preserving
the fixed-teacher rate.

\begin{corollary}[Two hidden layers]
\label{cor:two-layer-scale-refinement}
For $H=2$, the conclusions of
Theorem~\ref{thm:fixed-all-gamma-raw} hold under the weaker sample
condition
\begin{equation}
 n\geq C_2(d+1)^2
 \log^6\!\frac{2n(d+1)}{\lambda\epsilon\delta}
                   \lambda^{-9}\epsilon^{-4},
 \label{eq:two-layer-refined-budget}
\end{equation}
with the same restrictions $0<\lambda\leq\min\{\epsilon,\lambda_2\}$
and $\epsilon\leq1$.
Consequently one can compete, with vanishing additive excess, with
teachers whose width and weight caps are both at most $n^\alpha$
for every $\alpha<1/36$, simultaneously at every eligible learner
width and every $\gamma\geq1$.  At fixed weight cap, the sufficient
width range is $\alpha<1/18$.  For a fixed teacher the rate remains
$\widetilde O(n^{-1/13})$.
\end{corollary}

\begin{proof}
The deterministic approximation in
Lemma~\ref{lem:effective-approximation} has the sharper exact form
\[
 \|z_w-z_{\widetilde w}\|_\infty\leq8\tau\sqrt{q_N(w)}
 \qquad(H=2),
\]
where $\tau=\sigma\sqrt N$ is the actual normalized Gaussian scale.
Use the constant $\overline M_2=7T$ from the proof of
Theorem~\ref{thm:fixed-all-gamma-raw}, and set
\[
 \kappa=\frac{\sqrt{\lambda\epsilon}}
                    {96\sqrt{\overline M_2}}.
\]
The statistical calibrations and the bound $M_n\leq\overline M_2$
follow from~\eqref{eq:two-layer-refined-budget} exactly as in that proof.
For $\tau<\kappa$ and
$\widehat R_S(z_w)+\lambda q_N(w)\leq T$, the output error is at most
\[
 8\kappa\sqrt{T/\lambda}
 =\frac{\sqrt{T\epsilon}}{12\sqrt{\overline M_2}}
 \leq\sqrt T.
\]
The smooth-loss argument therefore still bounds the truncated raw risk
by $M_n$.  Its raw squared-loss transfer costs at most
\[
 48\sqrt{M_n}\tau\sqrt{q_N}
 \leq\frac\lambda{16}q_N+
                 \frac{9216M_n\tau^2}{\lambda}
 \leq\frac\lambda{16}q_N+\epsilon.
\]
The weighted posterior tail and the teacher comparison are unchanged.
This establishes the spectral part with this larger threshold.

Above the threshold, the common entropy cover costs a factor
$\kappa^{-2}=9216\overline M_2/(\lambda\epsilon)$.
The large-width sample requirement becomes
$C_2(d+1)^2\log^6(\cdot)\lambda^{-9}\epsilon^{-3}$,
while the small-width requirement becomes
$C_2(d+1)^2\log^6(\cdot)\lambda^{-9}\epsilon^{-4}$.
The latter dominates, proving~\eqref{eq:two-layer-refined-budget}
with the same simultaneous events.

For growing caps take $\lambda=n^{-\theta}$ and
$\epsilon=n^{-\eta}$ with
$4\alpha<\theta<1/9$, $0<\eta\leq\theta$, and
$9\theta+4\eta<1$.
These choices exist exactly in the stated sufficient range
$\alpha<1/36$.  For fixed weights replace $4\alpha$ by $2\alpha$.
Setting $\epsilon=\lambda$ gives sample power $13$.
\end{proof}


\section{Scope and remaining questions}
The fixed-prefactor results give polynomial oracle bounds simultaneously for every $\gamma\ge1$ and all learner widths above the comparison width, with the ordinary raw posterior mean. They require bounded inputs and labels, a specified competitor weight cap, and the stated likelihood calibration. They do not require target realizability or a correct observation-noise model.

Whether fixed-prefactor mean-field scaling attains the $\widetilde O(n^{-1/H})$ rate remains open. Our covariance cover proves a polynomial bound there, but its exponent worsens exponentially with hidden depth. Even at stronger scales, the unrestricted tanh bound has a sample exponent that worsens with depth. These results therefore do not give a polynomial jointly in teacher width and depth, nor a guarantee based only on parameter count.

The likelihood calibration is essential to the conclusion. Very tight conditioning can favor memorization of noise; excessive regularization can prevent approximation of a rapidly varying regression function. The positive results do not cover an arbitrary likelihood variance, efficient posterior computation, or equivalence with gradient descent.

\appendix
\section{Why width alone is insufficient}\label{width-only:sec:width-only-obstruction}\subsection{A finite-temperature obstruction}
Let $\pi$ be a prior on real-valued functions $z_w:[-1,1]\to\R$. For a sample $S=((x_i,y_i))_{i=1}^n$ with $|y_i|\le1$, write
\[
 \widehat R_S(z_w)=\frac1n\sum_{i=1}^n(z_w(x_i)-y_i)^2,
 \qquad
 \Pi_S(dw)=Z_S^{-1}e^{-\beta\widehat R_S(z_w)}\pi(dw),
 \quad 0<\beta<\infty.
\]
A Gaussian observation likelihood of variance $v>0$ has $\beta=n/(2v)$. The prior, its architecture, and $\beta$ are fixed before the data law below is chosen. They may depend on sample size.

Assume that each $z_w$ has a Lipschitz constant $\mathcal L(w)$ on $[-1,1]$ and that
\begin{equation}
 M=\sup_{|x|\le1}\E_\pi z_w(x)^2<\infty,
 \qquad A=\E_\pi\mathcal L(w)^2<\infty.
 \label{width-only:eq:moments}
\end{equation}
These conditions hold for every fixed finite feedforward tanh or ReLU architecture with Gaussian weights. No prior over network sizes is involved. Set
\[
 B_0=2(M+1),\qquad C=e^{\beta B_0/2}\sqrt A.
\]
For a prior whose output has mean zero, including an independent centered Gaussian readout, one may improve $B_0$ to $M+1$.

\begin{theorem}[A close-pair obstruction with fixed input covariance]
\label{width-only:thm:general}
For $0<\eta<1/2$, let $D_\eta$ put probability $1/4$ on each of
\[
 (\eta,1),\quad(-\eta,-1),\quad(1,0),\quad(-1,0).
\]
This law is noiseless and has input mean zero and variance $(1+\eta^2)/2\ge1/2$. For every sample $S$ of size $n$ from its support, the ordinary posterior mean $\bar z_S=\int z_w\,d\Pi_S$ satisfies
\begin{equation}
 R_{D_\eta}(\bar z_S)
 \ge\frac12(1-C\eta)_+^2.
 \label{width-only:eq:risk-lower}
\end{equation}
The same lower bound holds for the mean of clipped outputs, the clipped ordinary mean, and raw or clipped posterior-average risk. In particular, choosing $C\eta\le1/2$ gives a lower bound of $1/8$, uniformly over every possible sample.
\end{theorem}
\begin{proof}
For any bounded-label sample, Jensen and $(a-b)^2\le2a^2+2b^2$ give
\[
 Z_S\ge e^{-\beta\E_\pi\widehat R_S}\ge e^{-\beta B_0}.
\]
The likelihood is at most one. Thus, uniformly over all such samples,
\[
 \E_{\Pi_S}\mathcal L^2\le e^{\beta B_0}A=C^2.
\]
It follows from Cauchy--Schwarz that
\[
 |\bar z_S(\eta)-\bar z_S(-\eta)|
 \le2\eta\E_{\Pi_S}\mathcal L\le2C\eta.
\]
Write $u=\bar z_S(\eta)$ and $v=\bar z_S(-\eta)$. The two inner points alone contribute
\[
 \frac{(u-1)^2+(v+1)^2}{4}
 \ge\frac12\left(1-\frac{u-v}{2}\right)^2
 \ge\frac12(1-C\eta)_+^2
\]
to the risk. This proves~\eqref{width-only:eq:risk-lower}. Clipping is $1$-Lipschitz, so both clipped-mean conventions have the same pair-difference bound. Jensen then proves the posterior-average risk assertions. The moment assumptions and the evidence bound also ensure that the ordinary means and risks used here are finite.
\end{proof}

The bound tends to $1/2$ as $\eta\downarrow0$, although it is enough below to retain the fixed constant $1/8$. The argument does not depend on the sample containing an inner point: it is uniform over all samples, including those containing both oppositely labeled inner points many times.

\subsection{Exact small teachers}
The teacher class contains ordinary fully connected networks with zero values permitted for unused weights. Let $H\ge1$ be the number of hidden layers, followed by a linear readout. Only the first layer needs biases; internal biases may be fixed to zero.

\begin{lemma}[Two tanh chains]
\label{width-only:lem:tanh}
For every fixed $H\ge1$ and all sufficiently small $\eta>0$, a tanh network of hidden width two and depth $H+1$ fits $D_\eta$ exactly.
\end{lemma}
\begin{proof}
Let $F_H$ denote the $H$-fold composition of tanh and set
\[
 c_\eta=\frac{F_H(2/\eta)}{F_H(1/\eta)},\qquad
 a_\eta=F_H(2)-c_\eta F_H(1).
\]
The function $F_H$ is odd and strictly increasing. Moreover,
$F_H(t)\to F_{H-1}(1)>0$ as $t\to\infty$, with the convention $F_0(1)=1$. Consequently $c_\eta\to1$ and
$a_\eta\to F_H(2)-F_H(1)>0$. In particular $a_\eta>0$ for sufficiently small $\eta$.

Use two disjoint chains with first-layer weights $2/\eta$ and $1/\eta$, respectively, and internal weights one. The linear readout realizes
\[
 g_\eta(x)=\frac{F_H(2x/\eta)-c_\eta F_H(x/\eta)}{a_\eta}.
\]
By construction $g_\eta(\eta)=1$ and $g_\eta(1)=0$; oddness supplies the two negative inputs. Its readout coefficients remain bounded as $\eta\downarrow0$ at each fixed depth. Its first-layer weights diverge as $1/\eta$.
\end{proof}

\begin{lemma}[A ReLU tent]
\label{width-only:lem:relu}
For every $H\ge1$ and $0<\eta<1/2$, a ReLU network of hidden width four and depth $H+1$ fits $D_\eta$ exactly, with output magnitude at most one on the whole real line.
\end{lemma}
\begin{proof}
Write $r(t)=\max(t,0)$. The function
\[
 g_\eta(x)=\frac{-r(x+2\eta)+2r(x+\eta)-2r(x-\eta)+r(x-2\eta)}{\eta}
\]
is an odd piecewise-linear tent, equal to $\pm1$ at $\pm\eta$, zero outside $[-2\eta,2\eta]$, and bounded by one. Its four first-layer features are nonnegative and pass unchanged through identity internal matrices. The displayed coefficients give the readout. First-layer weights are one, bias magnitudes are at most $2\eta$, internal weights are zero or one, and readout magnitudes are at most $2/\eta$.
\end{proof}

\subsection{Gaussian priors and a width-only oracle}
\begin{lemma}[Finite Gaussian priors satisfy the required moments]
Every fixed finite fully connected tanh or ReLU network with Gaussian weights, a linear readout, and biases confined to the first layer satisfies~\eqref{width-only:eq:moments}. The conclusion allows any fixed positive Gaussian coordinate variances.
\end{lemma}
\begin{proof}
Both activations are $1$-Lipschitz and vanish at zero. If $W^1$ includes the bias column, then on $[-1,1]$ the augmented input has Euclidean norm at most $\sqrt2$. Put
\[
 T(w)=\|v\|_2\prod_{\ell=1}^H\|W^\ell\|_F.
\]
The output is at most $\sqrt2 T(w)$ in magnitude, and its input Lipschitz constant is at most $T(w)$. In a fixed finite architecture, $T(w)^2$ is bounded by a polynomial in the absolute values of finitely many Gaussian coordinates, so its expectation is finite. Independence is not needed for this integrability statement.
\end{proof}

\begin{corollary}[No width-only, distribution-uniform vanishing excess]
\label{width-only:cor:oracle}
Fix a hidden depth $H\ge1$. For each sample size $n$, choose any finite learner width $N_n$, a Gaussian weight prior with positive coordinate variances, and any finite positive inverse temperature $\beta_n$. Suppose $N_n\ge2$ for tanh, or $N_n\ge4$ for ReLU. There exist noiseless laws $D_n$ on $[-1,1]\times[-1,1]$, each with input variance at least $1/2$, and teachers of the corresponding fixed width such that
\[
 \inf_{g\text{ in the fixed-width teacher class}}R_{D_n}(g)=0,
 \qquad
 R_{D_n}(\bar z_{S,n})\ge\frac18
\]
for every sample $S$ of size $n$ from $D_n$. The other prediction conventions of Theorem~\ref{width-only:thm:general} satisfy the same lower bound.
\end{corollary}
\begin{proof}
For each prescribed learner and temperature, apply Theorem~\ref{width-only:thm:general} with $\eta_n>0$ small enough that $C_n\eta_n\le1/2$ and the corresponding exact-teacher lemma applies. The teacher may be zero-padded to any larger width. Both activations vanish at zero, so padding preserves its function.
\end{proof}

Thus no guarantee with a remainder tending to zero uniformly over input laws can depend only on teacher width when coefficient size is unrestricted. The counterexample already has fixed input dimension, bounded inputs and labels, no label noise, and input variance bounded away from zero. A coefficient, sensitivity, or approximation-scale restriction is needed for such a uniform statement. The result does not exclude an oracle inequality with a penalty depending on those additional resources.

\paragraph{Order of quantifiers.}
The prior and temperature are prescribed before selecting $\eta_n$. The result permits arbitrary deterministic schedules in sample size and learner width, including canonical Gaussian scaling regimes, provided every architecture and temperature is finite at each $n$. It does not exhibit one law on which every temperature fails. An adaptive rule may respond to the observed input separation by making the likelihood arbitrarily sharp; the proof does not bound such a rule. Nor does the theorem concern exact conditioning or a zero-temperature/strong-prior limit taken after fixing the data law.

\paragraph{Distribution-uniform versus pointwise learning.}
The laws $D_n$ and their exact teachers may vary with $n$. The statement refutes a uniform vanishing remainder over unrestricted-coefficient teachers; it does not refute consistency on one fixed law or one fixed teacher. The small separation supplies an unbounded sensitivity resource that is invisible to width and covariance alone.


\paragraph{Inputs on a unit sphere.}
The same obstruction holds with all inputs on the unit circle in $\mathbb R^2$. Replace the near inputs by $(\pm\eta,\sqrt{1-\eta^2})$ and the anchors by $(\pm1,0)$, retaining their labels and masses. The close-pair distance is still $2\eta$, the teachers use only the first coordinate, and the centered covariance is
\[
 \operatorname{diag}\bigl((1+\eta^2)/2,(1-\eta^2)/4\bigr).
\]
Its eigenvalues are at least $3/16$ when $\eta\le1/2$. The Gaussian moment and evidence proof applies on the bounded two-dimensional input set. Thus a unit-norm input assumption also does not replace control of teacher coefficients or sensitivity.


\begin{thebibliography}{9}
\bibitem{sst}
Nathan Srebro, Karthik Sridharan and Ambuj Tewari.
\emph{Smoothness, Low Noise and Fast Rates}.
Advances in Neural Information Processing Systems 23, 2010.
\url{https://www.cs.cornell.edu/~sridharan/smoothness.pdf}.
\bibitem{previous}
ChatGPT-6 Astra, Claude Opus 5.5, Kaarel H\"anni and Dmitry Vaintrob.
\emph{Sharper agnostic bounds for dense Gaussian Bayesian neural networks}.
Research manuscript, November 2026.
\url{https://github.com/kaarelh/AI-works/tree/main/neural-network-generalization/works-in-progress/agnostic-oracle-improvements}.
\end{thebibliography}
\end{document}
